\section{Open parameter statements II: assembly}\label{sec:pf:open0}

\providecommand{\openzeroMI}{\mathrm{MI}^{*}_{2}}
\providecommand{\openzeroHS}[1]{\|#1\|_{\mathrm{HS}}}
\providecommand{\openzerodev}{\operatorname{dev}}
\providecommand{\openzeromfL}{\operatorname{mfL}}
\providecommand{\openzerokfk}{\operatorname{kfk}}
\providecommand{\openzeroLk}{\mathrm{Lk}}
\providecommand{\openzeroEv}{\mathrm{Ev}}
\providecommand{\openzeroRow}{\sfname{Row}}
\providecommand{\openzeroLinked}{\sfname{LinkedAt}}
\providecommand{\openzeroAll}{\mathrm{All}}
\providecommand{\openzeroEvW}{\mathrm{EvW}}
\providecommand{\openzeroEvWC}{\mathrm{EvWC}}
\providecommand{\openzeroSC}{\mathrm{SC}}
\providecommand{\openzeroId}{\mathrm{Id}}
\providecommand{\openzeroDeg}{\mathrm{Deg}}
\providecommand{\openzeroNG}{\mathrm{NG}}
\providecommand{\openzeroGen}{\mathrm{Gen}}
\providecommand{\openzeroNZ}{\mathrm{NZ}}
\providecommand{\openzeropart}{\operatorname{part}}
\providecommand{\openzerovp}{\operatorname{vp}}
\providecommand{\openzerops}{\operatorname{ps}}
\providecommand{\openzeropa}{\operatorname{pa}}
\providecommand{\openzerobig}{\operatorname{big}}
\providecommand{\openzeroone}{\operatorname{one}}
\providecommand{\openzerogam}{\gamma}
\providecommand{\openzeroPg}{P_{\mathrm g}}
\providecommand{\openzeroGood}{\mathcal{G}}
\providecommand{\openzeroBuilt}{\mathcal{B}}
\providecommand{\openzeropdep}{p_{\mathrm{dep}}}
\providecommand{\openzeroDF}{\sfname{DefiniteSideCondition}}
\providecommand{\openzeroTrue}{E^{\mathrm{true}}}

This section proves the theorem of the formal development \cite{LeanSource2026} which supplies the twelve
\emph{open parameter statements} of \cref{thm:families:open}: from the package $P$ of classical results on unitary
groups it derives, for every proof $h_L$ of $\sfname{RecTExists}$ and every $c\in\R$, twelve linked rows for the
parameter $\openzeropdep(h_L,c)$. The multiplicity row among them is the subject of the previous section
(\cref{sec:pf:open0a}), which proves it under a \emph{diagonal row hypothesis}; here we derive this hypothesis from the
Siegel--Weil formula in the convergent range (\cref{lit:families:sw}), and we assemble the twelve statements.

\subsection{Introduction}\label{sec:pf:open0:intro}

\paragraph{The result.}
Each open parameter statement is a statement $\openzeroLinked(p,r)$ of the row language of \cref{def:families:rows}:
there is an auxiliary slot $R$ which, eventually and for large twisting data, contains every built point, such that the
row $r$ holds for the slot of good points of $R$ and the bundle linkage $\openzeroLk_p$ of $p$. Each row asserts that a
property of bundles holds at all bundles selected along the families of a class (eight of the twelve statements), or
eventually along each family (the goodness, multiplicity and $\kappa$-rows), or, for the definite-side row,
eventually and for all large constants $C$. The twelve properties are listed in \cref{pf:open0:tab}; the main theorem is
\cref{pf:open0:main}.

\paragraph{The argument.}
A bundle is linked to a point only through a \emph{bundle witness}, a record which identifies the real package of the
bundle with the real package of a twisted input constructed at the point. Three of the twelve statements are read off
the witness and its consequences directly (primitivity, decomposition points and the comparison); the geometric row
combines the comparison with a theorem on class-sum divisors proved elsewhere; two more statements need one further
input each (the join, from Harish-Chandra theory, and the $\lambda$-rows, from twisted multiplicity one). The goodness and
label rows come from Green's theorem on cuspidal characters and Zelevinsky's
theorem through the points of the construction; the $\kappa$-row is deduced from the $\lambda$-rows; the division row
is deduced from the multiplicity row and the law carried by the points of $\openzeropdep(h_L,c)$; and the definite-side
row is a height bound for column chains which holds at every linked bundle of every parameter.

The analytic content proved in full here is the \emph{diagonal row theorem} (\cref{pf:open0:diagthm}): for a word $u$ of
a datum, the normalised Gram operator $K_u=(c_0(u)^2)^{-1/2}P_u^*P_u$ of the period map of $u$ satisfies
\[
\openzeroHS{K_u-1}^2\ \le\ 4\bigl(e^{C\,m^5fL}+1\bigr)\frac{d_u^2}{D_*}+\bigl(4e^{-\frac12kf\log k}+2e^{-kf\log k}\bigr)\,d_u ,
\]
uniformly under the size conditions. The proof expands $K_u$ over the \emph{terms} of an orbital expansion (as in
Theorem~S, \cref{thm:theorem-S}, where two distinct words are compared). For equal words the identity terms sum to the
identity operator, the degenerate terms are negligible, and the remaining part $K_0-1$ is controlled by a Gram identity
over Levi types: its squared Hilbert--Schmidt norm is bounded by twice a ``price'' $\openzeroPg\le\frac65$ times the sum of the
pair sums of corank at least two (\emph{collisions}) and of corank one. The corank-one pairs are bounded through the
Siegel--Weil formula (generic pairs) and the non-generic row bound of \cref{prop:families:volratio}(iii).

\paragraph{What is proved in full; what is cited.}
We prove in full: the assembly of the twelve statements (\cref{sec:pf:open0:assembly}); the statements read off the
witness, at the level of their formal proofs (\cref{sec:pf:open0:witness}); the reduction of the goodness, label,
$\kappa$- and division rows to their inputs (\cref{sec:pf:open0:good,sec:pf:open0:mi}); the pointwise diagonal theorem,
with the abstract diagonal bound and all numerical lemmas, the column-dimension window, the decomposition of collisions
by corank and the deduction of the corank-one bounds from the two rows (\cref{sec:pf:open0:diag}); and the assembly of
the definite-side row (\cref{sec:pf:open0:df}); and, in \cref{sec:pf:open0:inputs}, the inputs of these arguments which belong
to the same part of the development: the ranks, the kernel ties and sizes and the kernel column theorem of the definite side, the type
labels and the label ties, the non-vanishing of isotypic components and the identification of $\kappa$, the division compatibility,
the identity cells of the orbital expansion, the pair amplitude bound, the amplitude, tie and size inputs of the collision bound and
the generic corank-one row. The deeper ingredients are stated precisely and cited; the largest of those which belong to the formal
region converted here are the kernel-half bound for column chains (about $1.2\cdot10^4$ Lean lines), the reduction modulo the maximal
compact subgroup behind the label ties (about $2.5\cdot10^3$), the plane mass bound behind the size hypothesis (about $1.8\cdot10^3$),
and the volume identity of the identity cells and the eigenvalue identity behind $\kappa$ (about $8\cdot10^2$ each).

\paragraph{Organisation.}
\Cref{sec:pf:open0:setting} recalls the row language, defines the twelve bundle properties and states the main
theorem; \cref{sec:pf:open0:assembly} proves it from the results of
\cref{sec:pf:open0:witness,sec:pf:open0:good,sec:pf:open0:mi,sec:pf:open0:diag,sec:pf:open0:df}, whose inputs are proved in
\cref{sec:pf:open0:inputs}. \Cref{sec:pf:open0:remarks} collects remarks on the formal statements.

%==============================================================================
\subsection{Rows, linkage and the twelve statements}\label{sec:pf:open0:setting}

\paragraph{Bundles and rows.}
We use the notions of \cref{def:families:closed,def:families:rows} in the precise form of
\cref{pf:rows:defrow,pf:rows:deflinked}. A \emph{bundle} $b$ over a datum $D$, an object $O$ over $D$ and $c\in\R$
consists of a reciprocity record $b.R$, real spectral ties $b.S$ over $b.R$, a row-side family $b.P$ determined by
$b.R$, a pair of natural numbers $(\Delta,h_0)$ such that $\Delta$ is the parameter $\Delta$ of the spectral data of $D$ and
$mh_0$ is the $\eta$-size of the size data of $O$, square-integrable columns
$b.C$, and a real theta package $b.B$ whose indefinite realisation we denote by $b.B.Q$. A \emph{linkage} is a relation
$\openzeroLk(K,\beta,c,b)$ between a context $K$, a bundle point $\beta$ of $K$ (with datum $\beta.D$ and object
$\beta.O$), a real number $c$ and a bundle $b$ over $(\beta.D,\beta.O,c)$; a \emph{slot} $Q$ is a predicate on the cores
of bundle points.

\begin{definition}[Rows and their targets]\label{pf:open0:def:rows}
Let $Q$ be a slot and $\openzeroLk$ a linkage.
\begin{enumerate}[label=(\alph*)]
\item For a property $E(\kappa,D,O,c,b)$ of bundles depending on a row parameter $\kappa$, $\openzeroEv_{Q,\openzeroLk}(E)$
means: eventually in the row parameters (for small exponent and staged large constants, as in \cref{pf:rows:defrow}),
every positive $\kappa$ admits $x_0$ such that for every context $K$ governed by $\kappa$ with scale $x\ge x_0$, every
bundle point $\beta$ of $K$ in $Q$, every $c$ and every bundle $b$ with $\openzeroLk(K,\beta,c,b)$, the property
$E(\kappa,\beta.D,\beta.O,c,b)$ holds.
\item For a property $H(\mathcal A,\sigma)$ of classes $\mathcal A$ with raw selections $\sigma$,
$\openzeroRow(Q,\openzeroLk,H)$ means: there is $E$ with $\openzeroEv_{Q,\openzeroLk}(E)$ such that $H(\mathcal A,\sigma)$
holds for every class $\mathcal A$ and every $E$-selection $\sigma$ over $\mathcal A$ (a selection whose bundles are
built for a row parameter linked to the family and satisfy $E$).
\item For a property $B$ of bundles: $\openzeroAll(B)(\mathcal A,\sigma)$ means $B(\sigma(\mathcal F,x,i))$ for all
$\mathcal F\in\mathcal A$, $x>x_1(\mathcal F)$ and $i\in I_x$; $\openzeroEvW(B)(\mathcal A,\sigma)$ means that, closedly on
the class $\mathcal W_3(\mathcal A)$ (\cref{def:families:closed}), every family $\mathcal F$ has an $x_0$ with
$B(\sigma(\mathcal F,x,i))$ for all $x\ge x_0$, $x>x_1(\mathcal F)$ and $i\in I_x$. For a family $B_{C,C_H}$ of properties,
$\openzeroEvWC(B)(\mathcal A,\sigma)$ means that, closedly on $\mathcal W_3(\mathcal A)$ with closing constant $C_H$, every
family has a $C_0\ge0$ such that for every $C\ge C_0$ there is an $x_0$ with $B_{C,C_H}(\sigma(\mathcal F,x,i))$ for all
$x\ge x_0$, $x>x_1(\mathcal F)$ and $i\in I_x$.
\item For a parameter $p$ with slot system $p.s$ and bundle linkage $\openzeroLk_p$, and a property $r(S,\openzeroLk)$ of
slots and linkages, $\openzeroLinked(p,r)$ means: there is a slot $R$ such that, eventually in the scale and for large
twisting data, every built core whose base satisfies the base condition $\mathrm{KP}_x$ of \cref{pf:rows:deflinked}
satisfies $R$, and $r(\openzeroGood_p\wedge R,\openzeroLk_p)$ holds, where $\openzeroGood_p$ is the slot of good points
of $p$.
\end{enumerate}
\formal{GlobalTheory.Fc.RowE}\ \formal{GlobalTheory.Fc.Row}\ \formal{GlobalTheory.Fc.AtAll}\
\formal{GlobalTheory.Fc.Ind.EvW}\ \formal{GlobalTheory.Fc.Ind.EvWC}\ \formal{GlobalTheory.Fc.RL.LinkedAt}
\end{definition}

We use the following formal properties of rows, proved elsewhere.
\begin{enumerate}[label=(R\arabic*)]
\item\label{pf:open0:R1} If $\openzeroEv_{Q,\openzeroLk}(\kappa\mapsto B)$, then $\openzeroRow(Q,\openzeroLk,\openzeroAll(B))$
(\lref{GlobalTheory.Fc.RowE.row}: take $E=B$).
\item\label{pf:open0:R2} If $B$ holds at every bundle $b$ with $\openzeroLk(K,\beta,c,b)$ for every context $K$ and
every $\beta$ in $Q$, then $\openzeroEv_{Q,\openzeroLk}(B)$ (\lref{GlobalTheory.Fc.RowE.of_all}); eventual statements
are closed under finite conjunctions and implications.
\item\label{pf:open0:R3} If $\openzeroRow(Q,\openzeroLk,H)$ and $H\Rightarrow H'$ on all $E$-selections, then
$\openzeroRow(Q,\openzeroLk,H')$ (\lref{GlobalTheory.Fc.Row.imp}).
\item\label{pf:open0:R4} $\openzeroRow(\openzeroGood_p,\openzeroLk_p,H)$ implies
$\openzeroLinked\bigl(p,(S,\openzeroLk)\mapsto\openzeroRow(S,\openzeroLk,H)\bigr)$
(\lref{GlobalTheory.Fc.RL.LinkedAt.of_row}: take $R$ the trivial slot).
\item\label{pf:open0:R5} $\openzeroAll(B)$ implies $\openzeroEvW(B)$ (\lref{GlobalTheory.Fc.FourBoundsAtBundles.evW_of_atAll}:
every family has $x_0=0$, closedly with $\eps_1=1$).
\item\label{pf:open0:R6} Two linked rows for the same parameter give a linked row for the conjunction of their
targets (\lref{GlobalTheory.Fc.RowL.and}, intersecting the auxiliary slots).
\end{enumerate}
A bundle selected by an $E$-selection is \emph{built}, i.e.\ equal to a bundle linked to a point of the slot; properties
which hold at all bundles with a witness therefore pass to selected bundles (\lref{GlobalTheory.Fc.JoinLinked.of_built_of_bWit}).

\paragraph{The parameter and the witnesses.}
A parameter $p$ for $h_L$ consists of a row position, a norm-choice extension $N$ (the extra data demanded at each point),
an $S$-extension and a geometric predicate; $p(N,S,\Gamma)$ denotes the parameter with the fixed position $R_X$. The
parameter of \cref{thm:families:open} is $\openzeropdep(h_L,c)=p(N_{\mathrm{dep}},\mathrm{triv},\mathrm{all})$, with the
trivial $S$-extension and the always-true geometric predicate, where at every point the extension $N_{\mathrm{dep}}$ asks
for a value of the \emph{law extension} $N_{\mathrm{law}}$ and a list pin at level $c$; a value of $N_{\mathrm{law}}$ is a
value of the default extension together with a proof of the \emph{law} $\mathrm{Law}(D)$: if every value
$B_u(j_N\varphi,j_{EE}x)$ of the complex pairings $B_u$ attached to the words of $D$ ($u$ a word, $\varphi\in N_u$, $x\in\mathit{RR}$) is the image of an element of
$\mathit{LL}$, then these values are given by the pairing $B_k$ of the divisor datum,
\begin{equation}\label{pf:open0:eq:law}
B_u(j_N\varphi,j_{EE}x)=\iota\bigl(B_k(u,\varphi,x)\bigr)\qquad\text{for all }u,\varphi,x .
\end{equation}
A \emph{weighted point} $q$ of $p.s$ in a context $K$ carries a datum $D_q$, an object $O_q$, a value of the extension and
further witness data; it determines a bundle point. A \emph{bundle witness} for $q$ and a bundle $b$ over $(D_q,O_q,c)$
consists of a twisted input $x$ over the extra datum of $O_q$, a twisted bundle $y$ over $x$, that is, a real input
$y.y$ for $(b.R,b.S)$ with a recursive input $y.r$ and intersection data, such that the real package of $b$ equals the
real package of $y.y$, together with compatibility data: the record of $b$ is compatible with the point, the spectral
ties of $b$ are the real spectral matching of the saturated spectral package of the point, the stage, weights and pins
of $y$ are those of the point, and the column characters of $b$ are compatible. The bundle linkage $\openzeroLk_p$
relates $(K,\beta,c,b)$ if $(\beta.D,\beta.O,b)=(D_q,O_q,b')$ for some weighted point $q$ of $p.s$ in $K$ and some bundle
$b'$ admitting a witness for $q$.

\begin{definition}[The twelve bundle properties]\label{pf:open0:def:twelve}
For a bundle $b$ over $(D,O,c)$ let $X_b$, $e_{V,b}$ and $s_{EE,b}$ be the coefficient complex, the row embeddings and the
conjugate-linear map of its indefinite realisation (\cref{sec:pf:open0a:setup}). The twelve properties are those of
\cref{pf:open0:tab}; $\openzeroMI$ is the multiplicity identity of \cref{pf:open0a:def:mistar}, and the other
properties are the following conditions on the components of $b$.
\begin{itemize}
\item $B_{\mathrm{prim}}$: the primitivity rows of the realisation of $b$ hold (Hecke rationality, unit and equivariance
conditions, block condition and $\sfname{RowTest}$) and the types of $b.P$ are separated.
\item $B_\lambda$: the $\lambda$-rows of $b.P$: there is a map from finite types to indices such that the isotypic
subspaces of each fibre are independent and span the image of the corresponding row space, the eigenvalues of $D$ agree
with those of $b.P$ at the relevant types, and the row counts add up to the dimensions.
\item $B_{\mathrm{div}}$, $B_{\mathrm{dc}}$: the division rows and the decomposition-point condition of $b.P$.
\item $B_{\mathrm{join}}$: a join of the realisation $b.B.Q$ with $b.R$ and $b.S$ exists.
\item $B_{\mathrm{cmp}}$: a \emph{comparison} of $b$ exists: a recursion of the realisation together with the
$\psi$-tie of the row-side family, the degree-three normalisation of the realisation and a normalisation compatible
with the recursion.
\item $B_{\mathrm{lab}}$: there are a Hida setting $X$ for $D$, a good operator $G$ on $X$ and a linear isometry $\iota$ from
the Hilbert space of the spectral data of $D$ onto the Poincar\'e space of $X$ such that the automorphic instance of
$X$ matches $D$ (ranks, degree, root discriminant, signs, Harish-Chandra parameters, discrete series and lowest
$K$-types, $g=g_0g_1$ and the formula for $D_*$), $\iota\circ s_{EE}$ is conjugate onto the conjugate of the image of $G$,
and the label conditions hold (the labels of all finite types are labels, the label tuples are good, and the base row of
every finite type lies in the type space of its label).
\item $B_{\mathrm{good}}$: every finite type of $D$ is \emph{good}: for each type prime $g_k$ ($k=0,1$) its label at $g_k$ is
a family of monic squarefree polynomials over $\F_{g_k}$ of degree $r$ with the separation, avoidance, non-splitting and
large-irreducible-factor conditions with respect to the characters $\chi^T_k$ and their inverses.
\item $B_{\mathrm{geo}}$: some comparison $t$ of $b$ has a part whose point has $\iota$-point summands, the class-sum pin,
the compatibility conditions and the additional point condition of the definite side.
\item $B_\kappa$: for every complex embedding $e$ of $\mathit{LL}$, the archimedean normalisation $\kappa_R(e)$ of the
divisor datum equals the $\kappa$ of the column side of $b.P$ at the place of $e$.
\item $B^{\mathrm{df}}_{C,C_H}$: the definite-side condition $\openzeroDF(D,O,c,b,C,C_H)$ of \cref{sec:pf:open0:df}.
\end{itemize}
\formal{GlobalTheory.Fc.S0.OpenParameterStatements}\ \formal{GlobalTheory.Fc.R1.PrimitivityBundleCondition}\
\formal{GlobalTheory.Fc.HasCompatibility}\ \formal{GlobalTheory.Fc.LabLinked.Ties}\ \formal{Formalism.BasicNotions.GoodFTAll}
\end{definition}

{\small
\begin{longtable}{>{\raggedright\arraybackslash}p{0.09\textwidth}>{\raggedright\arraybackslash}p{0.36\textwidth}>{\raggedright\arraybackslash}p{0.1\textwidth}>{\raggedright\arraybackslash}p{0.33\textwidth}}
\caption{The twelve open parameter statements $\openzeroLinked\bigl(p,(S,\openzeroLk)\mapsto\openzeroRow(S,\openzeroLk,H)\bigr)$.}\label{pf:open0:tab}\\
\toprule
field & bundle property & target $H$ & proved in\\
\midrule
\endfirsthead
\toprule
field & bundle property & target $H$ & proved in\\
\midrule
\endhead
\bottomrule
\endfoot
$\mathsf{hPA}$ & $B_{\mathrm{prim}}$: primitive rows, separated types & $\openzeroAll$ & \cref{pf:open0:hpa}\\
$\mathsf{hLamP}$ & $B_\lambda$: $\lambda$-rows & $\openzeroAll$ & \cref{pf:open0:lam}\\
$\mathsf{hDivG}$ & $B_{\mathrm{div}}$: division rows & $\openzeroAll$ & \cref{pf:open0:div}\\
$\mathsf{hDcPtG}$ & $B_{\mathrm{dc}}$: decomposition points & $\openzeroAll$ & \cref{pf:open0:dcpt}\\
$\mathsf{selJ}$ & $B_{\mathrm{join}}$: a join exists & $\openzeroAll$ & \cref{pf:open0:selj}\\
$\mathsf{tie}$ & $B_{\mathrm{cmp}}$: a comparison exists & $\openzeroAll$ & \cref{pf:open0:tie}\\
$\mathsf{lab}$ & $B_{\mathrm{lab}}$: Hida setting, good operator, isometry, label ties & $\openzeroAll$ & \cref{pf:open0:lab}\\
$\mathsf{good}$ & $B_{\mathrm{good}}$: all finite types good & $\openzeroEvW$ & \cref{pf:open0:good}\\
$\mathsf{geo}$ & $B_{\mathrm{geo}}$: geometric compatibility at a comparison & $\openzeroAll$ & \cref{pf:open0:geo}\\
$\mathsf{mi}$ & $\openzeroMI(X_b,e_{V,b},s_{EE,b})$ & $\openzeroEvW$ & \cref{pf:open0:mi}\\
$\mathsf{kap}$ & $B_\kappa$: $\kappa$-rows & $\openzeroEvW$ & \cref{pf:open0:kap}\\
$\mathsf{df}$ & $B^{\mathrm{df}}_{C,C_H}$: definite-side condition & $\openzeroEvWC$ & \cref{pf:open0:df}\\
\end{longtable}
}

\begin{theorem}[\cref{thm:families:open}]\label{pf:open0:main}
Assume the package $P$ of \cref{sec:families:theta}. For every proof $h_L$ of $\sfname{RecTExists}$ and every $c\in\R$,
the twelve statements of \cref{pf:open0:tab} hold for the parameter $\openzeropdep(h_L,c)$: for each row, with bundle
property $B$ and target $H$,
\[
\openzeroLinked\bigl(\openzeropdep(h_L,c),\ (S,\openzeroLk)\mapsto\openzeroRow(S,\openzeroLk,H(B))\bigr).
\]
\formallabel{GlobalTheory.T.open0_of_unitary_group_hypotheses}\
\formallabel{GlobalTheory.Fc.S0.Open0ProofU.open_orbit_of_unramified_exponents_alt}
\end{theorem}

The fields of $P$ which are used are: Harish-Chandra finiteness and rapid decay of cusp forms \cite{HarishChandra1968}, the
discrete series of $\Uni(p,q)$ and its covers and their existence \cite{HarishChandra1966}, Harish-Chandra modules and analyticity
\cite{HarishChandra1953}, Casselman--Wallach globalisation \cite{Casselman1989,Wallach1992}, trace class of cusp forms,
Dixmier--Malliavin and Flath factorisation \cite{DixmierMalliavin1978,Flath1979}, Weil representations of unitary dual pairs, Weil's
Siegel formula, the Siegel--Weil formula in the convergent range \cite{Siegel1935,Weil1965} (in the two formulations of the package;
the second is \cref{lit:families:sw}), the doubling method and the unramified doubling integral, the classification of discrete
series, strong approximation for $\mathrm{SU}$, Zelevinsky's polynomial ring theorem \cite{Zelevinsky1981}, the geometric lemma,
Green's theorem on cuspidal characters \cite{Green1955} and its version for regular characters, and the existence of ray class
fields. Howe's unramified computation \cite{Howe1979} is not a field of $P$: it is a theorem of the development
(\lref{GlobalTheory.Hk3.howeUnram_holds}).
The precise forms are those of \cref{sec:pf:open0a:setup} and \cite{LeanSource2026}.

%==============================================================================
\subsection{Proof of the main theorem}\label{sec:pf:open0:assembly}

Fix $h_L$ and $c$, and write $p=\openzeropdep(h_L,c)$, $\openzeroGood=\openzeroGood_p$, $\openzeroLk=\openzeroLk_p$. Two
facts about the extension $N_{\mathrm{dep}}$ are used repeatedly.
\begin{enumerate}[label=(L\arabic*)]
\item\label{pf:open0:L1} Every value of $N_{\mathrm{dep}}$ at a point yields a value of $N_{\mathrm{law}}$ at the same point
(its first component).
\item\label{pf:open0:L2} For every context $K$ and every weighted point $q$ of $p.s$ in $K$, the law $\mathrm{Law}(D_q)$ of
\eqref{pf:open0:eq:law} holds: it is a component of the value of $N_{\mathrm{law}}$ contained in the extension value of
$q$.
\end{enumerate}

\begin{proof}[Proof of \cref{pf:open0:main}]
The twelve statements are obtained as follows; each proposition cited is proved for the parameters $p(N,S,\Gamma)$ with
the indicated property of $N$, and $p=p(N_{\mathrm{dep}},\mathrm{triv},\mathrm{all})$.
\begin{enumerate}[label=\arabic*.]
\item $\mathsf{hPA}$ and $\mathsf{hDcPtG}$ hold for every parameter (\cref{pf:open0:hpa,pf:open0:dcpt}).
\item $\mathsf{tie}$ holds for every parameter (\cref{pf:open0:tie}); the formal statement used is made for parameters whose
extension yields the law extension, which $N_{\mathrm{dep}}$ does by \ref{pf:open0:L1}.
\item $\mathsf{selJ}$ follows from $\mathsf{tie}$ by \cref{pf:open0:selj}, using finiteness, rapid decay, the discrete
series cover, Harish-Chandra modules and Casselman--Wallach globalisation.
\item By \cref{pf:open0:mi}, the inputs of \cref{pf:open0a:main} together with the Siegel--Weil formula give two linked
rows: $\mathsf{mi}^{\mathrm{all}}$, with target $\openzeroAll(\openzeroMI)$, and $\mathsf{mi}$, with target
$\openzeroEvW(\openzeroMI)$; the latter is the multiplicity statement.
\item $\mathsf{hDivG}$ follows from \ref{pf:open0:L2} and $\mathsf{mi}^{\mathrm{all}}$ by \cref{pf:open0:div}.
\item $\mathsf{hLamP}$ holds for every parameter of the form $p(N,S,\Gamma)$ by \cref{pf:open0:lam}, and $\mathsf{kap}$ follows from
it by \cref{pf:open0:kap}.
\item Green's theorem in the form used by \cref{pf:open0:good} is deduced from the field on regular characters
(\lref{GlobalTheory.GreenBridge.green_of_regular}). \Cref{pf:open0:good} and \ref{pf:open0:L1} give two linked rows:
$\mathsf{good}^{\mathrm{all}}$, with target $\openzeroAll(B_{\mathrm{good}})$, and $\mathsf{good}$, with target
$\openzeroEvW(B_{\mathrm{good}})$.
\item $\mathsf{lab}$ follows from $\mathsf{good}^{\mathrm{all}}$ by \cref{pf:open0:lab}.
\item $\mathsf{geo}$ follows from $\mathsf{hDivG}$, $\mathsf{tie}$, \ref{pf:open0:L1} and the existence of ray class fields by
\cref{pf:open0:geo}.
\item $\mathsf{df}$ holds for every parameter (\cref{pf:open0:df}).
\end{enumerate}
The twelve statements are the fields of the conclusion.
\end{proof}

%==============================================================================
\subsection{Statements read off the bundle witness}\label{sec:pf:open0:witness}

In this subsection $p$ is an arbitrary parameter and $\openzeroLk=\openzeroLk_p$, unless stated otherwise. All proofs
have the same first step: if $\openzeroLk(K,\beta,c,b)$, there are a weighted point $q$ in $K$, a bundle $b'$ over
$(D_q,O_q,c)$ with a witness $w$, and an equality $(\beta.D,\beta.O,b)=(D_q,O_q,b')$; a property of bundles holds at $b$ as
soon as it holds at $b'$. We write $x$ for the twisted input of $w$ and $y$ for its twisted bundle; the real package of
$b'$ equals that of $y.y$, and in particular the row-side family of $b'$ is the \emph{division pair} $P_x$ of the adelic
forms, face and column data of $x$.

\begin{proposition}[Primitivity]\label{pf:open0:hpa}
Let $p$ be a parameter and $Q$ a slot. Then $\openzeroRow(Q,\openzeroLk_p,H)$ holds for the target
$H=\openzeroAll(B_{\mathrm{prim}})$; in particular $\openzeroLinked(p,\mathsf{hPA})$.
\formallabel{GlobalTheory.T.linkedAt_hPA_prmR}\ \formallabel{GlobalTheory.Fc.RL.PrimAt.linkedAt_hPA}\
\formallabel{GlobalTheory.Fc.RL.PrimAt.herm_wCell_last_mul_mulVec_eq_of_inWb}\ \formallabel{GlobalTheory.Fc.RL.PrimAt.of_wit_alt}
\end{proposition}

\begin{proof}
By \ref{pf:open0:R1} and \ref{pf:open0:R2} it suffices to show $B_{\mathrm{prim}}(b')$ for every bundle $b'$ with a witness
$w$. The intersection data of the twisted bundle $y$ contain the primitivity rows of the realisation of the real package
of $y.y$ and the separation of the types of its row-side family; since this real package equals that of $b'$, these
are the two conditions of $B_{\mathrm{prim}}(b')$. The linked row follows by \ref{pf:open0:R4}.
\end{proof}

\begin{lemma}[Rows at the companion]\label{pf:open0:companion}
For every twisted input $x$ (over an extra datum of a datum $D$ with gated ties and size data), the division pair $P_x$
satisfies: (a) the $\lambda$-condition $\mathcal L_3$ on the reading data of $x$ implies $B_\lambda(P_x)$; (b) the
decomposition-point rows of $P_x$ hold, and hence the decomposition-point condition $B_{\mathrm{dc}}(P_x)$; (c) the
division rows $\mathcal D_3(x)$ for the components of $x$ imply $B_{\mathrm{div}}(P_x)$.
\formallabel{GlobalTheory.Mem.R3.rowConditions_at_companion_from_readings}
\end{lemma}

\begin{proof}
(a) and (c) hold by the construction of $P_x$, whose $\lambda$-rows and division rows are by definition those of the
reading data, resp.\ of the components of $x$ (for (c) one passes from the division rows to their global form,
\lref{Formalism.L2.DIV.DivRows3.toG}). (b) The decomposition-point rows of $P_x$ are the point conditions of the
column data of $x$, and they imply the global condition (\lref{Formalism.L2.DIV.dcPtG_of}).
\end{proof}

\begin{proposition}[Decomposition points]\label{pf:open0:dcpt}
For every parameter $p$ and every slot $Q$, $\openzeroRow(Q,\openzeroLk_p,\openzeroAll(B_{\mathrm{dc}}))$; in particular
$\openzeroLinked(p,\mathsf{hDcPtG})$.
\formallabel{GlobalTheory.T.linkedAt_hDcPtG_prmR}\ \formallabel{GlobalTheory.Fc.RL.linkedAt_hDcPtG}\
\formallabel{GlobalTheory.Fc.RL.rowE_dcPt_linked_alt}
\end{proposition}

\begin{proof}
As for \cref{pf:open0:hpa}: at a bundle $b'$ with a witness, the row-side family is $P_x$, and
\cref{pf:open0:companion}(b) gives $B_{\mathrm{dc}}(P_x)$.
\end{proof}

\begin{proposition}[Comparisons]\label{pf:open0:tie}
For every parameter $p$ and every slot $Q$, $\openzeroEv_{Q,\openzeroLk_p}(B_{\mathrm{cmp}})$; in particular
$\openzeroLinked(p,\mathsf{tie})$. (The formal statement is made for parameters $p(N,S,\Gamma)$ such that every value of $N$
yields a value of $N_{\mathrm{law}}$, and carries seven classical results as hypotheses: discrete series and their covers,
Harish-Chandra modules, Casselman--Wallach, Weil representations, finiteness and rapid decay; none of these hypotheses is used.)
\formallabel{GlobalTheory.T.linkedAt_tie_of_nonempty_extensionWithLaw}\ \formallabel{GlobalTheory.Fc.RL.at_tie_alt}\
\formallabel{GlobalTheory.Fc.TieOfLink.hasTie_of_linked}\ \formallabel{GlobalTheory.Fc.TieOfLink.nonempty_tie_of_wit}\
\formallabel{GlobalTheory.Fc.TieOfLink.tie_of_rpk_alt}
\end{proposition}

\begin{proof}
The eventual statement is the conjunction of six eventual statements, five of which are identifications read off
witnesses at contexts whose scale is that of the point (lemmas of the formal proof, cited here):
\begin{enumerate}[label=(\roman*)]
\item the zero section of the scaled data of the witness is the canonical one for the scale of the context
(\lref{GlobalTheory.T.bWit_zsc_F0_heq_F0K_local_of_sc_eq});
\item the columns of the bundle are the column data of the point (\lref{GlobalTheory.T.bWit_C_eq_cf_of_sc_eq});
\item the place set of the witness is the image of the places of the point
(\lref{GlobalTheory.T.bWit_Pg_eq_image_plMap_of_sc_eq});
\item the theta-level characters of the realisation are conjugate to the level characters of the point
(\lref{GlobalTheory.T.theta_level_conjugate_character_alt});
\item the local row condition of the principal level frame holds with the level characters of the point
(\lref{GlobalTheory.T.local_level_character_pin_of_witness});
\end{enumerate}
and the sixth is the eventual statement that, for large twisting data, every stage tuple over a framed base satisfies the
host hypotheses of its frames (\lref{GlobalTheory.Mem.HostFrames.hostHyp_frames_ev}). At a bundle $b'$ with a
witness over a point in a context with the right scale, (ii)--(iv) and the host frames give the normalisation pins of the
realisation of $y.y$ (\lref{GlobalTheory.Mem.R3.realPins_realTOf_alt}), and (i) gives its degree-three normalisation
(\lref{GlobalTheory.Mem.R3.realPins3_realTOf_alt}). The recursive input $y.r$ gives a recursion of the realisation;
by \lref{GlobalTheory.T.recursion_local_row_condition_alt} it satisfies the $\psi$-tie of $P_x$ and, given the pins and
the local row condition (v), the pins are compatible with the recursion. These data form a comparison of $b'$, because
the real package of $b'$ equals that of $y.y$. Transporting along $(\beta.D,\beta.O,b)=(D_q,O_q,b')$ gives
$\openzeroEv_{Q,\openzeroLk_p}(B_{\mathrm{cmp}})$, and \ref{pf:open0:R1}, \ref{pf:open0:R4} give the linked row.
\end{proof}

\begin{proposition}[Joins]\label{pf:open0:selj}
Assume Harish-Chandra finiteness, rapid decay, the discrete series cover, Harish-Chandra modules and Casselman--Wallach
globalisation. For every parameter $p=p(N,S,\Gamma)$, $\openzeroLinked(p,\mathsf{tie})$ implies $\openzeroLinked(p,\mathsf{selJ})$.
(The formal statement also assumes $\openzeroLinked(p,\mathsf{hPA})$, $\openzeroLinked(p,\mathsf{hDcPtG})$, the existence of
Weil representations and of discrete series; these are not used.)
\formallabel{GlobalTheory.T.entries_combined_real_alt}\ \formallabel{GlobalTheory.Fc.JoinLinked.at_selJ_alt}\
\formallabel{GlobalTheory.Fc.JoinLinked.linked_atAll_of_imp}\ \formallabel{GlobalTheory.Fc.JoinLinked.selJ_of_bWit}
\end{proposition}

\begin{proof}
\emph{Step 1: passing an implication through linked rows.} Let $A,B$ be properties of bundles such that $A(b)\Rightarrow
B(b)$ for every bundle $b$ over a weighted point which admits a witness. Then $\openzeroLinked(p,\openzeroRow(\cdot,\cdot,
\openzeroAll(A)))$ implies $\openzeroLinked(p,\openzeroRow(\cdot,\cdot,\openzeroAll(B)))$: keep the auxiliary slot and the
property $E$; every bundle of an $E$-selection is built, hence equal to a bundle with a witness, so $A$ at the selected
bundles implies $B$ there.

\emph{Step 2: a witness and a comparison give a join.} Let $b$ be a bundle over a weighted point $q$ with a witness $w$
and a comparison $t$. The comparison provides the normalisation pins of the realisation of the real package of $b$,
hence of $y.y$. The witness provides a saturated spectral package $(s_P,s_S)$ equal to that of the point with $b.S$ the
real spectral matching of $s_S$; transferring the row placement of the spectral stage of the point along the stage move
of the witness (\lref{GlobalTheory.Mem.R3.StageMove.spectralStage_point_rowPlacement_transfer}) and the
identification of the columns of $b.R$ with those of the point gives the see-condition, the evaluation tie and the
join row condition for the Hida setting of $s_P$ at the stage of $y.y$. Under the five classical inputs, these
data and the intersection data of $y$ give a join of the realisation of $y.y$ with $(b.R,b.S)$
(\lref{GlobalTheory.Fc.JoinLinked.join_linked_realization_alt}, which builds the join from the join subspace of the
real data lattice). Since the real packages agree, this is a join for $b$. Step~1 with $A=B_{\mathrm{cmp}}$ and
$B=B_{\mathrm{join}}$ concludes.
\end{proof}

\begin{proposition}[$\lambda$-rows]\label{pf:open0:lam}
Assume the existence of Weil representations of unitary dual pairs and of discrete series, Harish-Chandra finiteness,
rapid decay, Harish-Chandra analyticity, the classification of discrete series, trace class of cusp forms,
Dixmier--Malliavin and Flath factorisation, Zelevinsky's theorem, the geometric lemma and Green's theorem. Then
$\openzeroLinked(p,\mathsf{hLamP})$ for every parameter $p=p(N,S,\Gamma)$.
\formallabel{GlobalTheory.T.real_group_inputs_combined_alt_alt}\ \formallabel{GlobalTheory.Fc.RL.at_row_hLamP_alt_alt}\
\formallabel{GlobalTheory.Fc.RL.LamLinked.isotRank_of_wit_alt}
\end{proposition}

\begin{proof}
Eventually in the row parameters, every bundle point $\beta$ admits a parameter choice for its datum and size data
(\lref{GlobalTheory.Fc.RL.LamLinked.evP_params}, from the numerical size rows of the scale). Let $b'$ be a bundle with
a witness over a weighted point $q$ whose datum admits a parameter choice. The real input $y.y$ provides the host
hypothesis of $D_q$. By twisted multiplicity one (\lref{GlobalTheory.T.twisted_multiplicity_one}, applied with the
analytic inputs and Green's theorem to the reading data of the point, its pins, its exact twists, the host hypothesis,
the parameter choice and the types of the point) the twisted multiplicities of the point are equal; with the good
stability of the type idempotents (from Green's theorem) and the cuspidal decomposition (from finiteness and rapid
decay) this gives the isotypic rank condition of $D_q$ (\lref{GlobalTheory.Mem.R3.isotRank_of_multEq}). Finite types
via the eigensystems of the full Hecke algebra (\lref{GlobalTheory.T.finiteTypes_via_full_heckeAlgebra_eigensystems})
then give the $\lambda$-condition $\mathcal L_3$ for the reading data of $x$, and \cref{pf:open0:companion}(a) gives
$B_\lambda(P_x)$, which is $B_\lambda(b')$. By \ref{pf:open0:R2} the property ``$B_\lambda$'' holds eventually at all linked
bundles, and \ref{pf:open0:R1}, \ref{pf:open0:R4} give the linked row.
\end{proof}

\begin{proposition}[Geometric compatibility]\label{pf:open0:geo}
Assume the existence of ray class fields, and let $p=p(N,S,\Gamma)$ with every value of $N$ yielding a value of
$N_{\mathrm{law}}$. Then $\openzeroLinked(p,\mathsf{hDivG})$ and $\openzeroLinked(p,\mathsf{tie})$ imply $\openzeroLinked(p,\mathsf{geo})$.
\formallabel{GlobalTheory.T.linkedAt_geometricWeakTies_of_tie}\
\formallabel{GlobalTheory.Fc.GeoZeroAt.geometricWeakTies_of_linked_ties}\
\formallabel{GlobalTheory.Fc.AllAtMember.forall_linked_bundles_and_alt}
\end{proposition}

\begin{proof}
Under these hypotheses the shared theorem \lref{GlobalTheory.T.class_sum_divisor_all_alt} gives the linked row
whose target is: at every selected bundle and for \emph{every} comparison $t$ of it, the part of $t$ has $\iota$-point
summands, its point has the class-sum pin, the compatibility conditions and the additional point condition of the definite side, and the datum
is tied to its size data. By \ref{pf:open0:R6} we may combine it with $\mathsf{tie}$; at a selected bundle, the row
$\mathsf{tie}$ provides a comparison $t$, and the first row applied to $t$ gives $B_{\mathrm{geo}}$. Conclude by
\ref{pf:open0:R3}, keeping the auxiliary slot.
\end{proof}

%==============================================================================
\subsection{Goodness, labels and the \texorpdfstring{$\kappa$}{kappa}-row}\label{sec:pf:open0:good}

\begin{proposition}[Goodness]\label{pf:open0:good}
Assume Green's theorem on cuspidal characters and Zelevinsky's theorem, and let $p=p(N,S,\Gamma)$ with every value of
$N$ yielding a value of $N_{\mathrm{law}}$. Then
$\openzeroLinked\bigl(p,\openzeroRow(\cdot,\cdot,\openzeroAll(B_{\mathrm{good}}))\bigr)$ and $\openzeroLinked(p,\mathsf{good})$.
\formallabel{GlobalTheory.T.linked_bundles_good_of_local_langlands_alt_alt}\
\formallabel{GlobalTheory.Fc.GoodLinked.row_good_alt_alt}\ \formallabel{GlobalTheory.Fc.GoodLinked.linkedAt_goodFTAll}\
\formallabel{GlobalTheory.Fc.GoodLinked.largeBasePropertyG_goodRec_alt_alt}\
\formallabel{GlobalTheory.Fc.GoodLinked.goodFTAll_at_members_of_depthZero}\
\formallabel{GlobalTheory.Fc.GoodLinked.row_goodRec_alt}\ \formallabel{GlobalTheory.Fc.LgModel.rowL_of_good}
\end{proposition}

\begin{proof}
Here, unlike the other statements, the auxiliary slot of $\openzeroLinked$ is not trivial. Let $\mathcal R$ be the slot of
cores whose datum has all its finite types good.

\emph{Step 1: members at layer points are good.} Eventually and for large twisting data, at every member of a layer point
of the construction the datum satisfies $B_{\mathrm{good}}$. Indeed, the coefficient duty pin at each member
(\lref{GlobalTheory.T.coefficient_duty_pin_at_each_member}) and the conjugacy of the type vectors
(\lref{GlobalTheory.TypeVecConj.typeVec}: for the row data of the point, every compatible choice of components of the
type representation admits a representation of the finite group $G_0$ with an invariant subspace, conjugate to the
components, containing a non-zero vector fixed by the type idempotent) are, eventually, the two hypotheses of the
non-Galois twist type theorem (\lref{UpperTheory.Lit.Dz.nonGaloisTwistType}), which under Green's and Zelevinsky's
theorems concludes $B_{\mathrm{good}}$.

\emph{Step 2: built points lie in $\mathcal R$.} Eventually and for large twisting data, every built core lies in $\mathcal
R$: a built core comes from a member at a layer point, with the norm-choice data of the point, a value of $N$ (hence of
$N_{\mathrm{law}}$), the reading structure and the pins of the point as the inputs of Step~1.

\emph{Step 3: the row.} At the slot $\openzeroGood\wedge\mathcal R$ the row $\openzeroRow(\cdot,\openzeroLk_p,
\openzeroAll(B_{\mathrm{good}}))$ holds trivially: by \ref{pf:open0:R2} every bundle linked to a point of the slot has a good
datum. A row at $\openzeroGood\wedge\mathcal R$ together with Step~2 gives the linked row with auxiliary slot $\mathcal R$
(\lref{GlobalTheory.Fc.LgModel.rowL_of_good} combines the large-base condition with the row). Finally
$\openzeroAll$ implies $\openzeroEvW$ by \ref{pf:open0:R5}, keeping the auxiliary slot.
\end{proof}

\begin{proposition}[Labels]\label{pf:open0:lab}
Assume Green's theorem and Zelevinsky's theorem, and let $p=p(N,S,\Gamma)$ with every value of $N$ yielding a value of
$N_{\mathrm{law}}$. Then $\openzeroLinked\bigl(p,\openzeroRow(\cdot,\cdot,\openzeroAll(B_{\mathrm{good}}))\bigr)$ implies
$\openzeroLinked(p,\mathsf{lab})$.
\formallabel{GlobalTheory.T.good_ft_entries_alt_alt}\ \formallabel{GlobalTheory.Fc.LabLinked.row_lab_of_layer_alt_alt}\
\formallabel{GlobalTheory.Fc.LabLinked.row_lab_of_good_alt_alt}\ \formallabel{GlobalTheory.Fc.LabLinked.ties_of_built_genPin}\
\formallabel{GlobalTheory.Fc.LabLinked.ties_of_genPin_goodFTAll'}\ \formallabel{GlobalTheory.Fc.LabLinked.ties_of_nAtData_real}
\end{proposition}

\begin{proof}
\emph{The generator pin.} For every weighted point $q$, the value of $N_{\mathrm{law}}$ contained in the extension value of
$q$ carries the identity $\zeta_k=\mathrm{gen}_k$ ($k=0,1$) between the generators $\zeta_k$ of the group data of the
norm-choice record and the generators of the label data at the type primes $g_k$.

\emph{Goodness implies label ties at a point.} Let $q$ be a weighted point with the generator pin, and assume
$B_{\mathrm{good}}(D_q)$. The saturation $s$ of $q$ provides a Hida setting $X=s.X$, a good operator $G=s.G$ and the isometry
$\iota$ of its spectral isomorphism. The row frame parts of the reading data of $q$
(\lref{GlobalTheory.Mem.R3.RowFrameParts.main}) and the type labels obtained from the generator pin under Green's and
Zelevinsky's theorems (\cref{pf:open0:typelab}) are the hypotheses of the label ties (\cref{pf:open0:irred}), which, given
$B_{\mathrm{good}}(D_q)$, yield the three conditions of $B_{\mathrm{lab}}$ for $(X,G,\iota)$.

\emph{The row.} Keep the auxiliary slot and the property $E$ of the given linked row. A bundle $b$ of an $E$-selection is
built, so $D=D_q$ for a weighted point $q$; the given row yields $B_{\mathrm{good}}(D)$ at $b$, and the previous step
yields $B_{\mathrm{lab}}(b)$. Conclude by \ref{pf:open0:R3}.
\end{proof}

\begin{proposition}[$\kappa$-rows]\label{pf:open0:kap}
Assume Harish-Chandra finiteness, rapid decay, the discrete series cover, Harish-Chandra modules, Casselman--Wallach,
Harish-Chandra analyticity, the classification of discrete series, trace class, Dixmier--Malliavin and Flath
factorisation, Zelevinsky's theorem, the geometric lemma and Green's theorem. For every parameter $p=p(N,S,\Gamma)$,
$\openzeroLinked(p,\mathsf{hLamP})$ implies $\openzeroLinked(p,\mathsf{kap})$.
\formallabel{GlobalTheory.T.harish_chandra_inputs_combined_alt_alt}\
\formallabel{GlobalTheory.Fc.KapLinked.row_kap_of_row_lam_alt_alt}
\end{proposition}

\begin{proof}
Keep the auxiliary slot and the property $E$ of $\mathsf{hLamP}$. We show $\openzeroAll(B_\kappa)$ for every $E$-selection
$\sigma$ satisfying $\openzeroAll(B_\lambda)$. The selected bundles are built for a row parameter, the data of a family
admit parameter choices (a field of the family), and by induction on builtness it suffices to show: if $b'$ has a witness
over a weighted point $q$ whose datum admits a parameter choice and $B_\lambda(b')$ holds, then $B_\kappa(b')$. The
row-side family of $b'$ is $P_x$. Twisted multiplicity one (\lref{GlobalTheory.T.twisted_multiplicity_one}) gives the
equality of twisted multiplicities for the reading data of $q$; from it and Green's theorem the isotypic components of
the rows are non-zero (\cref{pf:open0:isot}). The identification of $\kappa$ (\cref{pf:open0:kappa}), applied with the
$\kappa_R$-field of the object of the point, gives $B_\kappa(P_x)$. Finally $\openzeroAll$ implies $\openzeroEvW$
(\ref{pf:open0:R5}).
\end{proof}

%==============================================================================
\subsection{The multiplicity and division rows}\label{sec:pf:open0:mi}

\begin{proposition}[Multiplicity]\label{pf:open0:mi}
Assume \ref{pf:open0a:P1}--\ref{pf:open0a:P5}, \ref{pf:open0a:P7}--\ref{pf:open0a:P13}, Howe's unramified computation and
the Siegel--Weil formula of \cref{lit:families:sw}. Then for every $h_L$ and every parameter $p=p(N,S,\Gamma)$:
\[
\openzeroLinked\bigl(p,\openzeroRow(\cdot,\cdot,\openzeroAll(\openzeroMI))\bigr)\quad\text{and}\quad
\openzeroLinked\bigl(p,\openzeroRow(\cdot,\cdot,\openzeroEvW(\openzeroMI))\bigr),
\]
where $\openzeroMI$ stands for the property $b\mapsto\openzeroMI(X_b,e_{V,b},s_{EE,b})$.
\formallabel{GlobalTheory.T.main_theorem_published_inputs_third_alt_alt}\ \formallabel{GlobalTheory.miStarRow_alt_alt}
\end{proposition}

\begin{proof}
By \cref{pf:open0:diagthm} below, the Siegel--Weil formula implies the diagonal row hypothesis
$(\mathrm{DR}_{E_{\mathrm{diag}}})$ of \cref{pf:open0a:def:rowdiag}. Every weighted point $q$ of $p.s$ carries a reading
structure on its reading data (a component of its witness data). Hence \lref{GlobalTheory.T.main_theorem_published_inputs_variant_alt_alt}
(the main theorem of \cref{sec:pf:open0a}: under \ref{pf:open0a:P1}--\ref{pf:open0a:P13} and
$(\mathrm{DR}_{E_{\mathrm{diag}}})$, if every weighted point of $p$ admits a reading structure, then
$\openzeroMI(X_b,e_{V,b},s_{EE,b})$ holds eventually at the bundles linked to good points of $p$) gives
$\openzeroEv_{\openzeroGood_p,\openzeroLk_p}(\openzeroMI)$. By \ref{pf:open0:R1} this is a row with target
$\openzeroAll(\openzeroMI)$; by \ref{pf:open0:R3} and \ref{pf:open0:R5} also a row with target $\openzeroEvW(\openzeroMI)$; and
\ref{pf:open0:R4} turns both rows into linked rows.
\end{proof}

The division row rests on a vanishing statement for column pairings at the base tuple $\mu_0$. Recall from
\cref{pf:open0a:def:cx} that a coefficient complex $X$ has, for every tuple $t$, a $\mathcal K$-vector space $F_R(t)$ with an
$\iota_0$-semilinear map $j_R(t)$ to a Hilbert space $F_C(t)$, and linear maps $\Theta_t\colon F_C(t)\to E$ into a Hilbert space with
true part $\openzeroTrue$; the face provides integers $n_V(t)$, members $u$ with column pairings $\mathrm{Pt}(u,t)$, and for every
$t$ a matrix $\Pi(t)$ with $n_V(t)$ columns (the row projection of the first face datum); let $U(t)=\ker\Pi(t)\subseteq
\mathcal K^{n_V(t)}$. A \emph{coordinate datum} at the base level consists of $\mathcal K$-valued sesquilinear forms $h_t$ on $F_R(t)$,
subspaces $F^0_R(t)\subseteq F_R(t)$ (the good parts) and $\mathcal K$-linear coordinate maps $r_t\colon\mathcal K^{d(t)}\to F_R(t)$;
let $V^0_C(t)\subseteq F_C(t)$ be the complex span of $j_R(t)(F^0_R(t))$.

\begin{lemma}[Base vanishing]\label{pf:open0:btv}
Let $X$ be a coefficient complex with row embeddings $e_V$ and a conjugate-linear map $s_{EE}$, and a coordinate datum, such that
for every tuple $t$:
\begin{enumerate}[label=(\roman*)]
\item $\langle j_R(t)x,j_R(t)y\rangle=\iota_0(h_t(x,y))$ for $x,y\in F_R(t)$;
\item $r_t$ takes values in $F^0_R(t)$;
\item $\Theta_t(V^0_C(t))\subseteq\openzeroTrue$ and $\Theta_t^*(\openzeroTrue)\subseteq V^0_C(t)$;
\item the complex span of $j_R(t)(r_t(\mathcal K^{d(t)}))$ is $\operatorname{im}(\Theta_t^*)\cap V^0_C(t)$;
\item $e_V(t)F\in F^0_R(t)$ and $h_t\bigl(r_t(x),e_V(t)F\bigr)=h_t\bigl(r_t(x),r_t(\Pi(t)F)\bigr)$ for all $x$ and $F$.
\end{enumerate}
Then $\Theta_t\bigl(j_R(t)(e_V(t)F)\bigr)=\Theta_t\bigl(j_R(t)(r_t(\Pi(t)F))\bigr)$ for all $t$ and $F$. If moreover $(V,a,U,\zeta)$ are data
satisfying \ref{pf:open0a:M1}--\ref{pf:open0a:M3} and \ref{pf:open0a:M5} of \cref{pf:open0a:def:mistar}, then
$\mathrm{Pt}(u,\mu_0)(\varphi,F)=0$ for every member $u$, every $\varphi$ and every $F\in U(\mu_0)$.
\formallabel{Formalism.IndefR.Cx.CxData.theta_eV}\ \formallabel{Formalism.IndefR.Cx.CxData.base_tuple_vanishing_of_hypotheses}
\end{lemma}

\begin{proof}
Fix $t$ and $F$, put $q=r_t(\Pi(t)F)$ and $\delta=e_V(t)F-q$. By (v) and (ii), $\delta\in F^0_R(t)$ and $h_t(r_t(x),\delta)=0$ for all $x$.
Hence $j_R(t)\delta\in V^0_C(t)$, and $\Theta_t(j_R(t)\delta)\in\openzeroTrue$ by (iii). For $e\in\openzeroTrue$, the vector $\Theta_t^*e$ lies in
$\operatorname{im}(\Theta_t^*)\cap V^0_C(t)$ by (iii), hence in the span of the vectors $j_R(t)r_t(x)$ by (iv); and by (i)
$\langle j_R(t)r_t(x),j_R(t)\delta\rangle=\iota_0(h_t(r_t(x),\delta))=0$. Therefore
$\langle e,\Theta_t(j_R(t)\delta)\rangle=\langle\Theta_t^*e,j_R(t)\delta\rangle=0$ for all $e\in\openzeroTrue$. Taking
$e=\Theta_t(j_R(t)\delta)$ gives $\Theta_t(j_R(t)\delta)=0$, which is the first assertion since $\Theta_t\circ j_R(t)$ is additive.

Now let $F\in U(\mu_0)$, so that $\Pi(\mu_0)F=0$ and $r_{\mu_0}(\Pi(\mu_0)F)=0$. By \eqref{pf:open0a:eq:MI} with $t=\mu_0$,
$\zeta(u,\mu_0)=1$ (\ref{pf:open0a:M3}) and the first assertion,
\[
\iota_0\bigl(\mathrm{Pt}(u,\mu_0)(\varphi,F)\bigr)=a(u)\bigl\langle s_{EE}\bigl(V\Theta_{\mu_0}(j_R(\mu_0)(0))\bigr),
P_u(U(u)\iota_0(\varphi))\bigr\rangle=0 .
\]
The embedding $\iota_0$ is injective, so $\mathrm{Pt}(u,\mu_0)(\varphi,F)=0$. Of the hypotheses on $(V,a,U,\zeta)$ only
$\zeta(u,\mu_0)=1$ and \ref{pf:open0a:M5} are used.
\end{proof}

\begin{proposition}[Division rows]\label{pf:open0:div}
Assume the discrete series cover, Harish-Chandra modules, Casselman--Wallach, Harish-Chandra finiteness and rapid decay.
Let $p$ be a parameter such that $\mathrm{Law}(D_q)$ holds at every weighted point $q$ of $p.s$. Then
$\openzeroLinked\bigl(p,\openzeroRow(\cdot,\cdot,\openzeroAll(\openzeroMI))\bigr)$ implies $\openzeroLinked(p,\mathsf{hDivG})$.
(The formal statement also assumes the existence of Weil representations and of discrete series; they are not used.)
\formallabel{GlobalTheory.T.bk_law_entries_alt}\ \formallabel{GlobalTheory.Fc.DivLinked.at_hDivG_main_alt}\
\formallabel{GlobalTheory.Fc.DivLinked.hDivG_linked_of_mIstar2}\ \formallabel{GlobalTheory.Fc.DivLinked.hDivG_of_mIstar2}\
\formallabel{GlobalTheory.Fc.DivLinked.divisionRows_pairing_vanishes_of_rational_pairing_law}\
\formallabel{GlobalTheory.Div.pairing_definite_base_vanish_of_miStar2}
\end{proposition}

\begin{proof}
By Step~1 of the proof of \cref{pf:open0:selj} it suffices to show: if $b'$ has a witness over a weighted point $q$ and
$\openzeroMI(X_{b'},e_{V,b'},s_{EE,b'})$ holds, then $B_{\mathrm{div}}(b')$. Since the real package of $b'$ is that of
$y.y$, it suffices to treat the realisation of $y.y$ and the row-side family $P_x$.

\emph{Base vanishing.} The intersection data of $y$ give the first-level admissibility $\mathcal O_3(1)$ of the family, the separation of
types and the primitivity rows, and the pins of $x$ give the weak characterisation of the real datum; by the shared
lemma \lref{Formalism.IndefR.Links.complex_data_conditions_base_level_of_primitive_links} these imply the base-level
conditions of the coefficient complex (isometric theta maps and the three level conditions at the base level). The
definite coordinates lie in the good part, and the evaluation tie holds by $\mathcal O_3(1)$. Hence
\cref{pf:open0:btv}, applied to the data $(V,a,U,\zeta)$ of the multiplicity identity, gives the base vanishing
(the formal statement packaging these steps is \lref{GlobalTheory.Div.pairing_definite_base_vanish_of_miStar2}).

\emph{Division rows.} The pairings $B_u$ of $D_q$ take values in $\iota(\mathit{LL})$
(\lref{GlobalTheory.T.divisionPairing_rational_on_rationalPoints'}, from the five classical inputs, the pins of $x$ and the
intersection data of $y$); this is the hypothesis of the law $\mathrm{Law}(D_q)$, which therefore gives
\eqref{pf:open0:eq:law}. The division compatibility (\cref{pf:open0:divcompat}(b), with the group tie of the division data of
$x$) applies, the remaining hypotheses being components of $x$ and of the extra datum. Hence $\mathcal D_3(x)$ holds, and
\cref{pf:open0:companion}(c) gives $B_{\mathrm{div}}(P_x)$.
\end{proof}

%==============================================================================
\subsection{The diagonal row hypothesis from the Siegel--Weil formula}\label{sec:pf:open0:diag}

We use the notation of \cref{sec:pf:open0a:words}. For a datum $D$ and a word $u$, $C_u$ is the space of automorphic forms
of $u$, $P_u\colon C_u\to H$ its period map, $c_0(u)$ its normalising constant,
\[
K_u=(c_0(u)^2)^{-1/2}\,P_u^*P_u,\qquad d_0(u)=\dim C_u,\qquad \openzerodev(u)=\openzeroHS{K_u-1}^2 ;
\]
$D_*$ is the main term, $\openzeromfL(L)=m^5fL$ and $\openzerokfk=kf\log k$. The ranks satisfy $m\ge2$ and $k\ge m+3\ge5$; the
type primes are $g_0,g_1$. We write $\openzerogam_L(C)=\exp(C\,m^5fL)$.

\begin{theorem}[Diagonal row theorem]\label{pf:open0:diagthm}
Assume the Siegel--Weil formula in the convergent range (\cref{lit:families:sw}). Then the diagonal row hypothesis
$(\mathrm{DR}_{E_{\mathrm{diag}}})$ of \cref{pf:open0a:def:rowdiag} holds: for all reals $C_g,c_b,C_u,c_T,C_T$ with $c_b>0$
and $c_T>0$ there are $C\ge1$, $C'_0$ and $C_{l,0}$ such that for all $C'\ge C'_0$, $C_l\ge C_{l,0}$, every datum $D$
which is true in the strong sense, every $L\in\R$ and every word $u$ of $D$: if $D$ has size data $L_2$ with $L_2.L=L$,
$1\le C_\Gamma(L_2)$ and $8C_\Gamma(L_2)\le C'$, if the rank identity $d_0(w)=d_w^*(T_w^k-1)^f$ holds at every word $w$, and
if the standing size conditions $\openzeroSC(D;L,C',C_l,C_g,c_b,C_u,c_T,C_T)$ hold, then
\begin{equation}\label{pf:open0:eq:diag}
\openzerodev(u)\ \le\ 4\bigl(e^{C\,\openzeromfL(L)}+1\bigr)\frac{d_0(u)^2}{D_*}
+\bigl(4e^{-\openzerokfk/2}+2e^{-\openzerokfk}\bigr)\,d_0(u).
\end{equation}
\formallabel{GlobalTheory.T.diagonalRow_of_siegelWeil_unitary}\
\formallabel{GlobalTheory.Core.Rk.rowDiag_of_siegelWeilUnitary}
\end{theorem}

The bound \eqref{pf:open0:eq:diag} is the diagonal bound $\Delta_u(L,C)$ of \cref{pf:open0a:def:diag}. The main term $d_0(u)^2/D_*$
is the size of $\openzeroHS{K_u-1}^2$ for a random subspace of dimension $d_0(u)$ (\cref{rem:families:random}); the theorem
says that $K_u$ is as close to the identity as this, up to the factor $4(e^{C\openzeromfL(L)}+1)$ and a term which is small
compared with $d_0(u)$.

\subsubsection{Size conditions and term systems}

\begin{definition}[Standing hypotheses]\label{pf:open0:def:standing}
Let $D$ be a datum and write $b=e^L$. The standing size conditions $\openzeroSC(D;L,C',C_l,C_g,c_b,C_u,c_T,C_T)$ are the
conjunction of the ties of $D$, the mid-volume formula for the dimensions of the words, the size inequalities
$\mathrm{SI}$, the conditions $\mathrm{NW}(D)$ and the trace frame, where:
\begin{itemize}
\item $\mathrm{SI}$ consists of $1<L$ and a list $\mathrm{HL}'$ of inequalities, among which $C'm^5fL\le k$, $m^4\le f\le k$,
$(kb)^8\le g_0\le b^{C_g}$ and $(kb)^8\le g_1\le b^{C_g}$; the others bound $\log f$, the root discriminant, the
Harish-Chandra parameters and their gaps, the level primes and $\Delta$ in terms of $k,m,f,r,L$ and the constants;
\item $\mathrm{NW}(D)$ consists of interlacing conditions on the Harish-Chandra data, the splitting of the level primes
$T_w$ and of $g_0,g_1$ in the host, the congruences $T_w\equiv1\bmod(g_0g_1)^2$, a discriminant identity, $T_w\ne g_0,g_1$,
and the positivity $d_0(w)>0$ of the dimensions.
\end{itemize}
For size data $L_2$ of $D$ (with $L=L_2.L$), the \emph{standing hypotheses} $\mathrm{SH}=\mathrm{SH}(D,L_2,C',C_l,\dots)$ are
$8C_\Gamma(L_2)\le C'$ together with $\mathrm{HL}'$; the \emph{standing numerical data} $\mathrm{SN}(D,C_0)$ are $\mathrm{NW}(D)$
with the positivity of the dimensions replaced by the \emph{dimension window}
\begin{equation}\label{pf:open0:eq:window}
\Gamma_w^{-2}D_*/C_0\ \le\ d_0(w)\ \le\ C_0\,\Gamma_w^{-2}D_*\qquad\text{for all words $w$;}
\end{equation}
and $\mathrm{H}_0(C_0)$ is $\mathrm{SH}\wedge\mathrm{SN}(D,C_0)\wedge(1\le C_0)$.
\formal{GlobalTheory.Core.Rk.HD}\ \formal{GlobalTheory.Core.Deg.HypL}\ \formal{GlobalTheory.Core.Deg.DataNW}\
\formal{Formalism.MainNotions.S.StandingHypotheses}\ \formal{Formalism.MainNotions.S.StandingNumericalData}\
\formal{Formalism.MainNotions.S.Cnt.Hyp0}
\end{definition}

\paragraph{Term systems.}
We use the term systems of \cref{pf:thmS:terms} with their notation: a set of terms $t$ with operators $\mathrm{op}(t)$ in a
Hilbert space $H$ and vectors $\mathrm{vec}(t)$ in a Hilbert space $H_1$, rows, types, a weight $G(t)$, the predicates
``identity'' and ``generic'', a Gram matrix $\Gamma(i,j)$ on the types, a real $\Gamma_1$, a corank $s_0(t,t')$; the sets
$\openzeroId$, $\openzeroDeg=\{G=0,\neg\mathrm{isId}\}$, $\openzeroNG$, $\openzeroGen$, $\openzeroNZ=\{G\ne0\}$; the parts
$\openzeropart(S)=\sum_{t\in S}\mathrm{op}(t)$ and $K_0=\openzeropart(\openzeroId)+\openzeropart(\openzeroGen)+\openzeropart(\openzeroNG)$;
the set $P(S)$ of pairs in $S\times S$ with the same row, $\openzerobig(S)$ and $\openzeroone(S)$ the pairs of $P(S)$ with
$s_0\ge2$, resp.\ $s_0=1$; the pair sums $\openzerops(Q)=\Gamma_1\sum_Q\Re\langle\mathrm{vec}\,t,\mathrm{vec}\,t'\rangle$ and
$\openzeropa(Q)=\Gamma_1\sum_Q|\langle\mathrm{vec}\,t,\mathrm{vec}\,t'\rangle|$; and the conditions $\mathrm{GE}(S)$ (Gram
expansion), $\mathrm{GB}(\openzeroPg)$ (Gram bound), $\mathrm{RO}$ (row orthogonality) and $\mathrm{CP}$ (positive corank). In
addition, for $\mathbf 1\in H$, the condition $\mathrm{IO}(\mathbf 1)$ is $\openzeropart(\openzeroId)=\mathbf 1$.

For a datum $D$, a Hida setting $X$ for $D$, a full system $Y$ over $X$ and words $u,u'$, let $\mathcal T(u,u')$ be the term
system of the orbital expansion relating the column spaces of $u'$ and $u$ (\cref{sec:pf:thmS:abs}); here $H$ is the
Hilbert--Schmidt space of maps from the column space of $u'$ to that of $u$. The pair $(X,Y)$ is \emph{concretely tied} if the
level subgroups are normalised and neat, the column space of every word $w$ has dimension $d_0(w)$, the main term of $X$ is
$D_*$, Wu decay and integrability of the archimedean coefficient hold, the adelic measures, characters and local volumes
are compatible, $D$ is locally tied, and the trace frame holds. \emph{Concrete data} for $D$ are a Hida setting $X$, a full
system $Y$ with finite spaces, an isometric identification of the spectral data of $D$ with those of $Y$, and a concrete tie
of $(X,Y)$. For concretely tied $(X,Y)$ we write $d_u=d_0(u)$, $dd(u)=d_u^2/D_*$, $\mathbf 1_u\in H$ for the identity of the
column space of $u$, and
\[
\openzeroPg(Y)=\prod_v\Bigl(1+\frac3{|\F_{q_v}|}\Bigr)
\]
for the \emph{price} of $Y$, the product over the places $v$ of the typed reduction of $Y$, with residue fields $\F_{q_v}$.
\formal{Formalism.L2.SG.RowsConc.ConcTie}\ \formal{Formalism.L2.SG.ThmSOwn.ConcK}\ \formal{Formalism.L2.SG.RowsConc.PgC}

\begin{definition}[The four point bounds]\label{pf:open0:def:points}
Let $(X,Y)$ be concretely tied, $u$ a word, $\mathcal T=\mathcal T(u,u)$ and $\Gamma\in\R$.
\begin{itemize}
\item[(VG)] \emph{Degenerate bound}: the family $(\|\mathrm{op}(t)\|)_{t\in\openzeroDeg}$ is summable and
$\|\openzeropart(\openzeroDeg)\|^2\le d_u\,e^{-\openzerokfk}$.
\item[(CO)] \emph{Collision bound}: $\openzeropa(\openzerobig(\openzeroNZ))\le(g_0g_1)^{-kf/2}\,dd(u)+e^{-\openzerokfk/2}\,d_u$.
\item[(G1)$_\Gamma$] $\openzerops(\openzeroone(\openzeroGen))\le\Gamma\,dd(u)$;\qquad (N1)$_\Gamma$:
$\openzerops(\openzeroone(\openzeroNG))\le\Gamma\,dd(u)$.
\end{itemize}
For $c\in\R$ and a property $\mathcal Q$ of concretely tied pairs and words, $\mathrm{On}_c(\mathcal Q)$ means: for all
$C_g,c_b,C_u,c_T,C_T$ with $c_b,c_T>0$ there is a threshold $t$ such that $\mathcal Q$ holds for all $C',C_l\ge t$, every datum $D$
with size data $L_2$, all concrete data $(X,Y,\dots)$ for $D$ and every word $u$, provided $1\le C_\Gamma(L_2)$, the ties of $D$, the
mid-volume formula and $\mathrm{H}_0(c)$ hold and $\mathcal T(u,u)$ is absolutely convergent. The \emph{corank-one bounds}
$\mathrm{One}^{\mathrm{gen}}_c$, $\mathrm{One}^{\mathrm{ng}}_c$ mean: for all $C_g,c_b,C_u,c_T,C_T$ with $c_b,c_T>0$ there are $C\ge0$
and a threshold $t$ such that (G1)$_{\gamma_L(C)}$, resp.\ (N1)$_{\gamma_L(C)}$, with $L=L_2.L$, holds for all $C',C_l\ge t$, all
$D$ and $L_2$, all concretely tied $(X,Y)$ and all words $u$, provided the ties of $D$ and $\mathrm{H}_0(c)$ hold and
$\mathcal T(u,u)$ is absolutely convergent.
\formal{GlobalTheory.Core.Deg.VanishingGramTermsBound}\ \formal{GlobalTheory.Core.Deg.CollS}\
\formal{GlobalTheory.Core.Deg.On}\ \formal{GlobalTheory.Core.Diag.OneOn}
\end{definition}

\begin{theorem}[The four point bounds]\label{pf:open0:fourpoints}
\begin{enumerate}[label=(\alph*)]
\item For every $c\in\R$, $\mathrm{On}_c(\mathrm{VG})$ and $\mathrm{On}_c(\mathrm{CO})$ hold.
\item The Siegel--Weil formula of \cref{lit:families:sw} implies $\mathrm{One}^{\mathrm{gen}}_c$ and $\mathrm{One}^{\mathrm{ng}}_c$
for every $c\in\R$.
\end{enumerate}
\formallabel{GlobalTheory.T.degenerateTerms_and_collisionBounds_uniform_widths}\
\formallabel{GlobalTheory.Core.Deg.degenerate_and_collisions_uncond}\
\formallabel{GlobalTheory.T.corankOnePairsFamilyBounds_of_siegelWeil}\
\formallabel{GlobalTheory.Core.Diag.oneOn_singleton_of_siegelWeil}\
\formallabel{GlobalTheory.Core.Diag.oneOn_singleton_of_siegelWeil_of_stepChoice}
\end{theorem}

The proof is given in \cref{sec:pf:open0:fourpf}. We first show how the four bounds give the diagonal bound.

\subsubsection{The abstract diagonal bound}

\begin{lemma}\label{pf:open0:threehalves}
Let $W,W'$ be complex Hilbert spaces with $W$ finite-dimensional and $A,B\colon W\to W'$ linear. Then
$\openzeroHS{A+B}^2\le\frac32\openzeroHS{A}^2+3\openzeroHS{B}^2$. In particular, for every datum and word,
\[
\openzerodev(u)\le\tfrac32\,\openzeroHS{K_u-K_u^{\mathrm{deg}}-1}^2+3\,\openzeroHS{K_u^{\mathrm{deg}}}^2,
\]
where $K_u^{\mathrm{deg}}$ is the degenerate part of $K_u$ (the part of the term expansion of $K_u$ over the terms of weight $0$
which are not identity terms).
\formallabel{GlobalTheory.Core.Rk.hsNormSq_add_le_three_halves}\
\formallabel{GlobalTheory.Core.Rk.hsDev_le_split_three_halves}\
\formallabel{GlobalTheory.T.gramPoincareHilbertSchmidtDeviation_le_split}
\end{lemma}

\begin{proof}
Let $(e_i)$ be an orthonormal basis of $W$, $a_i=\|Ae_i\|$, $b_i=\|Be_i\|$. Since $(a_i-2b_i)^2\ge0$, $2a_ib_i\le\frac12a_i^2+2b_i^2$,
hence $\|(A+B)e_i\|^2\le(a_i+b_i)^2\le\frac32a_i^2+3b_i^2$; sum over $i$. For the second statement write
$K_u-1=(K_u-K_u^{\mathrm{deg}}-1)+K_u^{\mathrm{deg}}$.
\end{proof}

\begin{lemma}[Abstract diagonal bound]\label{pf:open0:abstract}
Let $\mathcal T$ be an absolutely convergent term system with $H_1$ complete, let $\mathbf 1\in H$, and let $\openzeroPg,B,A_1,A_2$ be
reals with $\openzeroPg\ge0$. Assume $\mathrm{GE}(\openzeroGen)$, $\mathrm{GE}(\openzeroNG)$, $\mathrm{GB}(\openzeroPg)$,
$\mathrm{RO}$, $\mathrm{CP}$, $\mathrm{IO}(\mathbf 1)$, and
\[
\openzeropa(\openzerobig(\openzeroNZ))\le B,\qquad\openzerops(\openzeroone(\openzeroGen))\le A_1,\qquad
\openzerops(\openzeroone(\openzeroNG))\le A_2 .
\]
Then $\|K_0-\mathbf 1\|^2\le2\,\openzeroPg\,(B+A_1+A_2)$.
\formallabel{GlobalTheory.diag_abstract_collisions_once}\
\formallabel{GlobalTheory.T.abstractDiagonalBound_collisions_counted_once}
\end{lemma}

\begin{proof}
We use the lemmas on term systems proved with Theorem~S: for absolutely summable families the pair expansion
$\Phi(S)=\openzerops(P(S))$ with its splitting by corank (\lref{Formalism.FourierDecayAlt.P.GrowthPolynomialOperator.t_PairSplit}),
the monotonicity $\openzeropa(Q)\le\openzeropa(Q')$ for $Q\subseteq Q'$ and $|\openzerops(Q)|\le\openzeropa(Q)$
(\lref{Formalism.FourierDecayAlt.P.GrowthPolynomialOperator.t_SubFamily}), and the price of the types: under
$\mathrm{GE}(S)$ and $\mathrm{GB}(\openzeroPg)$, $\|\openzeropart(S)\|^2\le\openzeroPg\,\Phi(S)$
(\lref{Formalism.FourierDecayAlt.P.GrowthPolynomialOperator.t_PriceOfIota}), where $\Phi(S)=\Gamma_1\sum_i\|\openzerovp(S,i)\|^2$.

\emph{Step 1: splitting by corank.} By the pair expansion (with $\mathrm{RO}$ and $\mathrm{CP}$),
$\Phi(\openzeroNG)=\openzerops(\openzerobig(\openzeroNG))+\openzerops(\openzeroone(\openzeroNG))$. Applying the pair expansion to
the term system obtained from $\mathcal T$ by declaring every pair scalar (which changes neither $\Phi$ nor $\openzerobig$
nor $\openzeroone$, empties the non-scalar pairs and makes the scalar pairs of corank one equal to $\openzeroone$) gives
$\Phi(\openzeroGen)=\openzerops(\openzerobig(\openzeroGen))+\openzerops(\openzeroone(\openzeroGen))$.

\emph{Step 2: collisions are counted once.} The family $|\langle\mathrm{vec}\,t,\mathrm{vec}\,t'\rangle|$ is summable over pairs,
since it is bounded by $\|\mathrm{vec}\,t\|\,\|\mathrm{vec}\,t'\|$ and $\sum_t\|\mathrm{vec}\,t\|<\infty$. The sets
$\openzerobig(\openzeroGen)$ and $\openzerobig(\openzeroNG)$ are disjoint (a term is either generic or not), and both are contained in
$\openzerobig(\openzeroNZ)$ (generic and non-generic terms have non-zero weight). Hence
\[
\openzeropa(\openzerobig(\openzeroGen))+\openzeropa(\openzerobig(\openzeroNG))
=\openzeropa(\openzerobig(\openzeroGen)\cup\openzerobig(\openzeroNG))\le\openzeropa(\openzerobig(\openzeroNZ))\le B .
\]

\emph{Step 3: the two parts.} By the price of the types, Step~1, $\openzerops\le\openzeropa$ and the hypotheses,
\[
\|\openzeropart(\openzeroGen)\|^2\le\openzeroPg\bigl(\openzeropa(\openzerobig(\openzeroGen))+A_1\bigr),\qquad
\|\openzeropart(\openzeroNG)\|^2\le\openzeroPg\bigl(\openzeropa(\openzerobig(\openzeroNG))+A_2\bigr).
\]

\emph{Step 4.} By $\mathrm{IO}(\mathbf 1)$, $K_0-\mathbf 1=\openzeropart(\openzeroGen)+\openzeropart(\openzeroNG)$, and
$\|x+y\|^2\le2\|x\|^2+2\|y\|^2$. With Steps~2 and~3 and $\openzeroPg\ge0$,
$\|K_0-\mathbf 1\|^2\le2\openzeroPg\bigl(\openzeropa(\openzerobig(\openzeroGen))+\openzeropa(\openzerobig(\openzeroNG))+A_1+A_2\bigr)
\le2\openzeroPg(B+A_1+A_2)$.
\end{proof}

\begin{lemma}[The price]\label{pf:open0:price}
Let $D$ be a datum with size data $L_2$ and $X,Y$ a Hida setting and a full system for $D$, and assume the hypotheses
$\mathrm{H}_0(C_0)$ together with the parameter conditions of the $K_2$-theorems (in particular $(kb)^8\le g_0,g_1$, $f\le k$
and $b\ge2$). Then every residue field of the typed reduction of $Y$ has at least $3$ elements, and $\openzeroPg(Y)\le\frac65$.
\formallabel{GlobalTheory.T.equalIdentificationPairPrice_le_six_fifths}\ \formallabel{GlobalTheory.Core.Deg.pgC_le_six_fifths}\
\formallabel{GlobalTheory.Core.Deg.one_factor_sharp}\ \formallabel{GlobalTheory.Core.Deg.exp_pair_le_six_fifths}\
\formallabel{GlobalTheory.Core.Deg.exp_le_inv_one_sub_pg}\ \formallabel{Formalism.L2.SG.RowsConc.card_ge_three}
\end{lemma}

\begin{proof}
The places of the typed reduction lie above $g_0$ or $g_1$, and the residue fields there are $\F_{g_0}$, resp.\ $\F_{g_1}$,
since $g_0,g_1$ split completely; precisely, by the shared lemma \lref{Formalism.L2.SG.Pf.SatGram.tt_eq'} (from the
split-prime localisation and the reduction structure) and since $[F:\Q]=f$,
\[
\openzeroPg(Y)=\Bigl(1+\frac3{g_0}\Bigr)^f\Bigl(1+\frac3{g_1}\Bigr)^f .
\]
Since $k\ge5$ and $b>1$ we have $kb\ge5$ and $g_j\ge(kb)^8\ge3$, which is the first assertion. Let $g\in\{g_0,g_1\}$. As
$b\ge2$, $b^8\ge256$, and $k\le k^8$, so $256k\le(kb)^8\le g$; with $f\le k$ this gives $3f/g\le3/256$, and
$1+3/g\le e^{3/g}$ gives $(1+3/g)^f\le e^{3f/g}\le e^{3/256}$. Hence $\openzeroPg(Y)\le e^{3/128}$. Finally $e^x(1-x)\le1$ for
$x<1$ (as $1-x\le e^{-x}$), so $e^{3/128}\le128/125\le6/5$.
\end{proof}

\begin{lemma}[Numerics]\label{pf:open0:num}
\begin{enumerate}[label=(\alph*)]
\item If $x\ge2$, then $e^{-x}\le\frac25e^{-x/2}$.
\item Let $H,N_1,N_2,\openzeroPg,s,dd,d,E_2,E,\Gamma_1,\Gamma_2,\Gamma$ be reals with $H\le\frac32N_1+3N_2$,
$N_1\le2\openzeroPg(s\,dd+E_2d+\Gamma_1dd+\Gamma_2dd)$, $N_2\le dE$, $\openzeroPg\le\frac65$, $0\le s\le1$, $dd,d,E_2,\Gamma_1,\Gamma_2\ge0$,
$E\le\frac25E_2$ and $\Gamma_1+\Gamma_2\le\Gamma$. Then $H\le4\bigl((\Gamma+1)dd+dE_2\bigr)+2dE$.
\item If $y\ge1$, then $e^{C_1y}+e^{C_2y}\le e^{(\max(C_1,C_2)+1)y}$ for all real $C_1,C_2$.
\end{enumerate}
\formallabel{GlobalTheory.Core.Diag.exp_neg_le_two_fifths}\ \formallabel{GlobalTheory.Core.Diag.diag_arith}\
\formallabel{GlobalTheory.Core.Rk.exp_add_exp_le_exp_max_succ}
\end{lemma}

\begin{proof}
(a) $e^{-x}=e^{-x/2}\cdot e^{-x/2}$ and $e^{-x/2}\le e^{-1}<\frac25$ since $e>2.5$. (b) The bracket in the bound for $N_1$ is
non-negative; since $\openzeroPg\le\frac65$, $s\,dd\le dd$ and $(\Gamma_1+\Gamma_2)dd\le\Gamma dd$,
$\frac32N_1\le\frac{18}5\bigl((\Gamma+1)dd+dE_2\bigr)$; and $3N_2\le3dE=2dE+dE\le2dE+\frac25dE_2$. Adding,
$H\le\frac{18}5(\Gamma+1)dd+4dE_2+2dE$, which is at most the claimed bound. (c) With $M=\max(C_1,C_2)$, both terms are at
most $e^{My}$, and $e^y\ge1+y\ge2$, so the sum is at most $2e^{My}\le e^ye^{My}$.
\end{proof}

\begin{lemma}[Identity cells]\label{pf:open0:identity}
For every datum $D$, every concretely tied pair $(X,Y)$ and every word $u$, the identity part of $\mathcal T(u,u)$ is the
identity: $\mathrm{IO}(\mathbf 1_u)$ holds.
\formallabel{Formalism.L2.SG.RowsConc.idOneC_holds}\ \formallabel{Formalism.L2.SG.RowsConc.idOneC_of_full}\
\formallabel{Formalism.L2.SG.RowsConc.idOne_of_cell}
\end{lemma}

\begin{proof}
The concrete tie gives the normalisation of the level subgroups and Wu decay, hence folding at $u$
(\lref{Formalism.L2.SG.Closed.FoldOKC_outright}); fundamental sets exist and the orbital integrals over them are
integrable (\lref{Formalism.L2.SG.Closed.FundExistsC_holds} and \lref{Formalism.L2.SG.Closed.OOIntAllC_outright}). Under these
conditions \cref{pf:open0:idcells} applies: the operators of the
identity terms of the orbital expansion for the pair $(u,u)$, computed on a fundamental set, have unconditional sum the
identity of the column space of $u$. The embedding into the Hilbert--Schmidt space $H$ is linear and continuous, so the
identity part $\openzeropart(\openzeroId)$ of $\mathcal T(u,u)$ is $\mathbf 1_u$.
\end{proof}

\begin{proposition}[The pointwise diagonal theorem]\label{pf:open0:pointwise}
Let $D$ be a datum with size data $L_2$ and concrete data $(X,Y,\dots)$, let $C',C_l\ge29$, $C_0\in\R$, and assume
$\mathrm{H}_0(C_0)$ for $(D,L_2,C',C_l,C_g,c_b,C_u,c_T,C_T)$. Let $\Gamma_1,\Gamma_2\ge0$ with $\Gamma_1+\Gamma_2\le\Gamma$, and let $u$ be a
word such that (VG), (CO), (G1)$_{\Gamma_1}$ and (N1)$_{\Gamma_2}$ hold for $\mathcal T(u,u)$. Then
\[
\openzerodev(u)\le4\Bigl((\Gamma+1)\frac{d_0(u)^2}{D_*}+d_0(u)\,e^{-\openzerokfk/2}\Bigr)+2\,d_0(u)\,e^{-\openzerokfk}.
\]
\formallabel{GlobalTheory.T.diagonalTheorem_of_fourPointBounds}\ \formallabel{GlobalTheory.Core.Diag.sDiagAt_of_points}\
\formallabel{GlobalTheory.T.kernel_expansion_absolute_convergence_equal_words}\ \formallabel{GlobalTheory.Core.Diag.kExp_absConv_tied}\
\formallabel{Formalism.L2.SG.RowsConc.dd_nonneg}
\end{proposition}

\begin{proof}
\emph{Step 1: the expansion.} Folding, the Poincar\'e space, fundamental sets, orbital integrability, independence of the
fundamental set, well-definedness of the amplitudes, absolute convergence and the kernel expansion (eight shared
statements, combined by \lref{Formalism.L2.SG.ThmSOwn.kExpTied_of_full} and \lref{Formalism.L2.SG.RowsConc.absConvTied_of_full})
show that for concretely tied pairs the term system $\mathcal T=\mathcal T(u,u)$ is absolutely convergent and that $K_u$ is the
sum of its term operators. With the identification of spectral data, this gives
\[
\openzeroHS{K_u-K_u^{\mathrm{deg}}-1}^2=\|K_0-\mathbf 1_u\|^2,\qquad\openzeroHS{K_u^{\mathrm{deg}}}^2=\|\openzeropart(\openzeroDeg)\|^2
\]
(\lref{Formalism.L2.SG.ThmSOwn.hsNormSq_diag} and \lref{Formalism.L2.SG.ThmSOwn.hsNormSq_Kpp}; the first also uses that identity
terms have weight $0$, \lref{Formalism.L2.SG.Pf.SatGram.idDegenerate_tO}). By \cref{pf:open0:threehalves},
\begin{equation}\label{pf:open0:eq:split}
\openzerodev(u)\le\tfrac32\|K_0-\mathbf 1_u\|^2+3\|\openzeropart(\openzeroDeg)\|^2 .
\end{equation}

\emph{Step 2: the Gram data.} Since $C'\ge1$, $\mathrm{H}_0(C_0)$ implies the hypotheses of the $K_2$-theorems
(\lref{Formalism.MainNotions.NonGenericAndKTwoTheorems.Hyp_of}), so $\openzeroPg(Y)\le\frac65$ and all residue fields have at
least $3$ elements (\cref{pf:open0:price}). The Gram bound $\mathrm{GB}(\openzeroPg(Y))$ for $\mathcal T$ follows from the bound for
Gram matrices of character sums over $\GL_k(\F_q)$ (\lref{Formalism.L2.SG.Closed.TermsGGBound_holds}), the non-triviality of the
additive characters of $Y$ and the bound $|\F_{q_v}|\ge3$ (\lref{Formalism.L2.SG.TermOp.ggBound}). The Gram expansions
$\mathrm{GE}(\openzeroGen)$ and $\mathrm{GE}(\openzeroNG)$ hold for concretely tied pairs
(\lref{Formalism.L2.SG.RowsConc.normExpSatTied_holds}), as do $\mathrm{RO}$ (\lref{Formalism.L2.SG.TermOp.rowOrth}),
$\mathrm{CP}$ (\lref{Formalism.L2.SG.KappaOffDiagonal.corankPosC_holds}) and $\mathrm{IO}(\mathbf 1_u)$ (\cref{pf:open0:identity}).
Let $s=(g_0g_1)^{-kf/2}$ and $E_2=e^{-\openzerokfk/2}$. \Cref{pf:open0:abstract}, with $B=s\,dd(u)+E_2d_u$ from (CO),
$A_1=\Gamma_1dd(u)$ and $A_2=\Gamma_2dd(u)$, gives
\[
\|K_0-\mathbf 1_u\|^2\le2\openzeroPg(Y)\bigl(s\,dd(u)+E_2d_u+\Gamma_1dd(u)+\Gamma_2dd(u)\bigr).
\]

\emph{Step 3: numerics.} By the concrete tie, $d_u=d_0(u)$ and $dd(u)=d_0(u)^2/D_*\ge0$. Since $m\ge2$ and $k\ge m+3$, $k\ge5$;
since $m^4\le f$, $f\ge1$; as $k\ge5>e$, $\log k\ge1$, so $\openzerokfk\ge5\ge2$ and \cref{pf:open0:num}(a) gives
$E=e^{-\openzerokfk}\le\frac25E_2$. As $g_0,g_1$ are primes, $\sqrt{g_0g_1}\ge1$ and $0\le s\le1$. By (VG),
$\|\openzeropart(\openzeroDeg)\|^2\le d_uE$. \Cref{pf:open0:num}(b) applied to \eqref{pf:open0:eq:split} and the bounds of
Step~2 gives the assertion.
\end{proof}

\subsubsection{The dimension window and concrete data}

\begin{lemma}[Dimension window]\label{pf:open0:window}
\begin{enumerate}[label=(\alph*)]
\item Let $S>0$, $T\ge2$ and $1\le f\le k$ be integers ($f,k$) and reals ($S,T$) with $S\le T\le(1+\frac1{fk})S$. Then
$S^{fk}/2\le(T^k-1)^f\le3S^{fk}$.
\item If $a,G>0$ and $n\ge1$, then $\bigl(a^{1/n}G^{-2/n}\bigr)^n=a/G^2$.
\item If a datum $D$ satisfies $\mathrm{NW}(D)$ and the rank identity $d_0(w)=d_w^*(T_w^k-1)^f$ at every word $w$, then it
satisfies $\mathrm{SN}(D,3)$.
\end{enumerate}
\formallabel{GlobalTheory.Diag.PG.DuWinThree.window}\ \formallabel{GlobalTheory.Diag.PG.DuWinThree.star_pow}\
\formallabel{GlobalTheory.Diag.PG.duWin_three_of_count}\
\formallabel{GlobalTheory.Core.Diag.standingNumericalData_three_of_dataNW_rankMu}\
\formallabel{GlobalTheory.T.columnDimension_window_three_of_columnCount}
\end{lemma}

\begin{proof}
(a) Since $k-1<2^{k-1}$, $2k\le2^k$; hence $T^k\ge2^k\ge2k\ge2f$. Write $(T^k-1)^f=T^{fk}(1-T^{-k})^f$; by Bernoulli's
inequality $(1-T^{-k})^f\ge1-fT^{-k}\ge\frac12$, so $(T^k-1)^f\ge\frac12T^{fk}\ge\frac12S^{fk}$. Conversely $(T^k-1)^f\le
T^{fk}\le(1+\frac1{fk})^{fk}S^{fk}\le eS^{fk}\le3S^{fk}$, using $1+x\le e^x$. (b) is the rule $(x^\alpha)^\beta=x^{\alpha\beta}$ for $x>0$.
(c) Only the dimension window \eqref{pf:open0:eq:window} with $C_0=3$ has to be shown. Fix $w$ and write $T=T_w$, $\Gamma=\Gamma_w>1$,
$d^*=d^*_w\ge1$ and $D_*(w)>0$ for the main term of the member of $w$, which equals $D_*$. The member provides the window
$T_w^*\le T\le(1+\frac1{fk})T_w^*$ of \eqref{pf:open0a:eq:Tstar} for $T_w^*=(D_*(w)/d^*)^{1/(fk)}\Gamma^{-2/(fk)}$, and by (b),
$(T_w^*)^{fk}=c:=D_*(w)/(d^*\Gamma^2)$. Since $T\ge2$ is prime, $1\le f\le k$ ($f\ge m\ge2$), (a) gives $c/2\le(T^k-1)^f\le3c$. Multiplying by
$d^*$ and using the rank identity, $\Gamma^{-2}D_*/2\le d_0(w)\le3\Gamma^{-2}D_*$, which implies \eqref{pf:open0:eq:window} with
$C_0=3$.
\end{proof}

\begin{lemma}[Concrete data]\label{pf:open0:concrete}
Assume Wu decay in the form of the development (for every datum with the gap hypothesis and every Hida setting $X$ for it,
the Wu decay estimate holds and the archimedean coefficient of $X$ is integrable). Let $D$ be a datum with size data $L_2$,
spectral ties $S$ whose full system has finite spaces, local ties and the trace frame, and assume
$\mathrm{SH}(D,L_2,C',C_l,\dots)$ with $C',C_l\ge29$. Then concrete data for $D$ exist.
\formallabel{GlobalTheory.T.nonempty_concreteRowData_of_finitenessBundle}\ \formallabel{GlobalTheory.Core.Diag.concK_of_finX}
\end{lemma}

\begin{proof}
The standing hypotheses with $C_l\ge16$ and $C'\ge29$ imply the gap hypothesis (\lref{Formalism.L2.SG.ThmSOwn.gap_hypothesis_of_hypotheses}),
so Wu decay and integrability hold for the Hida setting $S.X$. The level subgroups of $S.X$ are normalised
(\lref{Formalism.L2.SG.Closed.level_subgroup_properties}) and hence neat at every word. Together with the finiteness,
the local ties and the trace frame these are the components of a concrete tie of $(S.X,S.Y)$, and with the spectral
identification of $S$ they form concrete data (\lref{Formalism.L2.SG.ThmSOwn.concKOf}).
\end{proof}

\subsubsection{Proof of the diagonal row theorem}

\begin{proposition}[Reduction to the four point bounds]\label{pf:open0:reduction}
Assume Wu decay in the form of \cref{pf:open0:concrete}, $\mathrm{On}_3(\mathrm{VG})$, $\mathrm{On}_3(\mathrm{CO})$,
$\mathrm{One}^{\mathrm{gen}}_3$ and $\mathrm{One}^{\mathrm{ng}}_3$. Then the conclusion of \cref{pf:open0:diagthm} holds.
\formallabel{GlobalTheory.Core.Rk.diagonalRow_bound_from_windowWidthThree}
\end{proposition}

\begin{proof}
Let $C_g,c_b,C_u,c_T,C_T$ with
$c_b,c_T>0$. The corank-one bounds give $C_1\ge0$, $C_2\ge0$ and thresholds $t_1,t_2$; the bounds on the width $3$ give
thresholds $t_G,t_C$. Put
\[
C=\max(C_1,C_2)+1\ \ (\ge1),\qquad C'_0=C_{l,0}=\max(t_1,t_2,t_G,t_C,29).
\]
Let $C'\ge C'_0$, $C_l\ge C_{l,0}$, $D$ true in the strong sense, $L\in\R$ and $u$ a word, with $E_{\mathrm{diag}}(D,L,C')$ and
$\openzeroSC$. The first gives size data $L_2$ with $L_2.L=L$, $8C_\Gamma(L_2)\le C'$, $1\le C_\Gamma(L_2)$, and the rank identity at all
words; the second gives the ties of $D$, the mid-volume formula, $\mathrm{SI}$, $\mathrm{NW}(D)$ and the trace frame. Hence
$\mathrm{SH}$ holds, $\mathrm{SN}(D,3)$ holds by \cref{pf:open0:window}(c), and so $\mathrm{H}_0(3)$. A world of $D$ provides spectral
ties $S$ with finite spaces and the local ties of $D$; by \cref{pf:open0:concrete} there are concrete data $(X,Y,\dots)$ for $D$, and
$\mathcal T(u,u)$ is absolutely convergent (Step~1 of the proof of \cref{pf:open0:pointwise}). Since $C',C_l$ exceed the four
thresholds, the four point bounds hold at $u$: (VG), (CO), (G1)$_{\Gamma_1}$ and (N1)$_{\Gamma_2}$ with $\Gamma_i=\gamma_L(C_i)\ge0$.
Put $y=m^5fL$; since $m\ge2$, $f\ge m^4\ge1$ and $L>1$, $y\ge1$, and \cref{pf:open0:num}(c) gives $\Gamma_1+\Gamma_2\le\gamma_L(C)=:\Gamma$.
\Cref{pf:open0:pointwise} (applicable since $C',C_l\ge29$) gives
\[
\openzerodev(u)\le4\Bigl((e^{C\openzeromfL(L)}+1)\frac{d_0(u)^2}{D_*}+d_0(u)e^{-\openzerokfk/2}\Bigr)+2d_0(u)e^{-\openzerokfk},
\]
which is \eqref{pf:open0:eq:diag}.
\end{proof}

\begin{proof}[Proof of \cref{pf:open0:diagthm}]
Wu decay holds (\lref{GlobalTheory.wuDecayOfGap}). By \cref{pf:open0:fourpoints} with $c=3$, the Siegel--Weil formula gives
the four point bounds on the width $3$, and \cref{pf:open0:reduction} concludes.
\end{proof}

\subsubsection{Proof of the four point bounds}\label{sec:pf:open0:fourpf}

\paragraph{The degenerate bound.}
For concretely tied $(X,Y)$ and a word $u$, the \emph{pair amplitude bound} is the conclusion of \cref{pf:open0:pairsamp}: the
family $(\|\mathrm{vec}\,t\|\,\|\mathrm{vec}\,t'\|)$ over the pairs of degenerate terms of $\mathcal T(u,u)$ in the same row is summable and
$\openzeroPg(Y)\,\Gamma_1$ times its sum is at most $d_u\,e^{-\openzerokfk}$; it is the hypothesis of the abstract degenerate bound
\lref{Formalism.L2.SG.RowsOther.pairsDeg_of_amp}.

\begin{proposition}\label{pf:open0:deg}
$\mathrm{On}_c(\mathrm{VG})$ holds for every $c\in\R$.
\formallabel{GlobalTheory.Core.Deg.vanishingGramBoundWindow_of_termPairTypeBound}\
\formallabel{GlobalTheory.Core.Deg.vanishingGramBound_uniformWindow_and_weightData}\
\formallabel{GlobalTheory.T.deg_termSumBoundOfPairsAmplitude}\ \formallabel{GlobalTheory.Core.Deg.termSumBound_of_pairsAmplitude}
\end{proposition}

\begin{proof}
\emph{Step 1: amplitudes give the degenerate bound.} Let $(X,Y)$ be concretely tied, $u$ a word with $\mathcal T=\mathcal T(u,u)$
absolutely convergent, and assume the $K_2$-hypotheses and the pair amplitude bound. The degenerate terms are summable
(absolute convergence). The Gram expansion holds on the degenerate part (\lref{Formalism.L2.SG.Pf.normExpDegC_of_leaves},
from the finiteness of Levi quotients and the summability of orbital integrals), and $\mathrm{GB}(\openzeroPg(Y))$,
$\mathrm{RO}$, $\mathrm{CP}$ hold as in Step~2 of \cref{pf:open0:pointwise}. The pair amplitude bound gives the degenerate pair bound
(\lref{Formalism.L2.SG.RowsOther.pairsDeg_of_amp}), and the abstract degenerate bound
(\lref{Formalism.L2.SG.RowsOther.kppDiag_of_pairs}, which combines the price of the types, the Gram sums, the pair
expansion and the monotonicity of pair sums as in \cref{pf:open0:abstract}) gives $\|\openzeropart(\openzeroDeg)\|^2\le
d_ue^{-\openzerokfk}$.

\emph{Step 2: the type bound gives the amplitudes.} The type bound of the development
(\lref{GlobalTheory.rowsPos_typeP}: for all constants there are $C_H,C_h,C_p\ge0$ such that for every $C_s\ge C_p+4$ there is
$C_\circ$ such that for every $C\ge C_\circ$ the type condition $\mathrm{Ty}(C_H,C_s,C,C_h,C_p)$ holds at all concretely tied
pairs under the $K_2$-hypotheses for $C',C_l$ beyond thresholds) is applied with $C_s=C_p+4$ and $C=\max(C_\circ,2(C_H+1))$. Under
the type condition the pair amplitude bound holds for $C_l$ beyond a threshold (\cref{pf:open0:pairsamp}). Taking for $C'$ the maximum of the thresholds and $1$ and for $C_l$
the maximum of the two thresholds, the $K_2$-hypotheses follow from $\mathrm{H}_0(c)$ (\lref{Formalism.MainNotions.NonGenericAndKTwoTheorems.Hyp_of}),
and Step~1 applies. This proves $\mathrm{VG}$ uniformly in the width; restricting to the width $c$ gives $\mathrm{On}_c(\mathrm{VG})$.
\end{proof}

\paragraph{The collision bound.}
For $2\le s\le k$ let $\openzerobig_s\subseteq P(\openzeroNZ)$ be the set of pairs of corank $s$.

\begin{lemma}[Decomposition by corank]\label{pf:open0:classes}
Let $\mathcal T$ be absolutely convergent, $k\in\N$, and $B\colon\N\to\R$. If $s_0(t,t')\le k$ for all $(t,t')\in\openzerobig(\openzeroNZ)$
and $\Gamma_1\sum_{(t,t')\in\openzerobig_s}\|\mathrm{vec}\,t\|\,\|\mathrm{vec}\,t'\|\le B(s)$ for $2\le s\le k$, then
$\openzeropa(\openzerobig(\openzeroNZ))\le\sum_{s=2}^kB(s)$.
\formallabel{GlobalTheory.T.pairAbs_corank_ge_two_le_sum_corankClasses}\ \formallabel{GlobalTheory.Core.Deg.pairAbs_big_le_sum_classes}
\end{lemma}

\begin{proof}
By the corank bound, $\openzerobig(\openzeroNZ)$ is the disjoint union of the sets $\openzerobig_s$, $2\le s\le k$. By Cauchy--Schwarz
$|\langle\mathrm{vec}\,t,\mathrm{vec}\,t'\rangle|\le\|\mathrm{vec}\,t\|\,\|\mathrm{vec}\,t'\|$, and the latter family is summable over all pairs
(absolute convergence). Since $\Gamma_1>0$, summing over the finitely many classes gives the claim.
\end{proof}

\begin{lemma}[Corank sums]\label{pf:open0:corank}
For reals $C,b,g,v,\Gamma$, naturals $k,m,f,s$ and a function $\Xi\colon\N\to\R$ put $H=b^{Cm^2f}$ and
\begin{align*}
N_s&=(2/g)^{(s-1)kf}(kv)^{kf}\,\Xi(s)\,C^{kf}\,H^{C(s+m)m^2}(Csf)^{6sf}\,\Gamma,\\
G_s&=(2/g)^{(s-2)kf}(kv)^{kf}\,\Xi(s)\,C^{kf}\,H^{C(s+m)m^2}(Csf)^{6sf}\,\Gamma\,e^{Cm^6f\log b}\,(C/k)^{kf}.
\end{align*}
Let $\mathrm{PC}(C,K_0)$ be the parameter conditions $4\le m\le f\le k$, $b\ge2$, $(kb)^8\le g$, $K_0m^5f\log b\le k$, $1\le v\le b$ and
$1\le N_T\le b^{Cmf}$ (for an auxiliary real $N_T$; the formal conditions also contain $1\le\Gamma'\le e^{Cm^5f\log b}$ for an auxiliary
$\Gamma'=1$, which holds automatically), and let $\mathrm{XH}(C)$ be the size hypothesis: $\Xi(s)\ge0$ and
$(kv)^{kf}\Xi(s)\le H^C\,C^{skf}$ for $2\le s\le k-2$, and $0\le\Xi(s)\le(Ck)^{kf}(H^{Cm}N_T)^{Ck}$ for $s\in\{k-1,k\}$. For every $C\ge1$
there is $K_0$ such that, under $\mathrm{PC}(C,K_0)$ and $\mathrm{XH}(C)$, for every $\Gamma$ with $0\le\Gamma\le e^{k/8}$:
\[
\sum_{s=2}^kN_s\le g^{-kf/2},\qquad\sum_{s=2}^kG_s\le e^{-\frac12kf\log k}.
\]
\formallabel{GlobalTheory.T.corankSums_bound_with_expLoss}\ \formallabel{GlobalTheory.Core.Deg.corankSums_expLoss}\
\formallabel{GlobalTheory.Core.Deg.gam_absorb}\ \formallabel{GlobalTheory.Core.Deg.ndRow_mono}\
\formallabel{GlobalTheory.Core.Deg.degRow_mono}\ \formallabel{GlobalTheory.Core.Deg.params_mono}\
\formallabel{GlobalTheory.Core.Deg.xiHyp_mono}
\end{lemma}

\begin{proof}
We use the two summation bounds of the development for the row functions with $\Gamma=1$
(\lref{Formalism.HermitianLatticeCounting.P.sumNondegenerateRowsLe} and \lref{Formalism.HermitianLatticeCounting.P.sumDegenerateRowsLe}):
for every $C^\sharp\ge1$ there is $K$ such that, under $\mathrm{PC}(C^\sharp,K)$, the additional condition $1\le\Gamma\le
e^{C^\sharp m^5f\log b}$ for $\Gamma=1$, and $\mathrm{XH}(C^\sharp)$, the sums of $N_s$ and of $G_s$ over $2\le s\le k$, formed with
$(C^\sharp,\Gamma=1)$, are at most $g^{-kf/2}$, resp.\ $e^{-\frac12kf\log k}$. Let $C\ge1$, $C^\sharp=eC$, and let $K_1,K_2$ be the
constants of the two bounds for $C^\sharp$; put $K_0=\max(K_1,K_2)$.

\emph{The hypotheses pass from $C$ to $C^\sharp$.} $\mathrm{PC}(C,K_0)$ implies $\mathrm{PC}(C^\sharp,K_i)$: the condition
$K_im^5f\log b\le k$ follows from $K_i\le K_0$ and $\log b\ge0$, and $N_T\le b^{Cmf}\le b^{C^\sharp mf}$ since $b\ge1$; moreover
$1\le e^{C^\sharp m^5f\log b}$. Since $b\ge1$, $H=H(C)=b^{Cm^2f}\ge1$ is increasing in $C$; with $N_T\ge1$ this gives
$\mathrm{XH}(C)\Rightarrow\mathrm{XH}(C^\sharp)$, each upper bound being increasing in $C$.

\emph{Absorbing $\Gamma$.} Since $f\ge1$, $k/8\le kf$, so $C^{kf}\Gamma\le C^{kf}e^{kf}=(eC)^{kf}$ for $0\le\Gamma\le e^{k/8}$ and $C\ge0$.
Comparing factor by factor ($H(C)^{C(s+m)m^2}\le H(C^\sharp)^{C^\sharp(s+m)m^2}$ as $H(C)\ge1$; $(Csf)^{6sf}\le(C^\sharp sf)^{6sf}$; and, for
$G_s$, $e^{Cm^6f\log b}\le e^{C^\sharp m^6f\log b}$ and $(C/k)^{kf}\le(C^\sharp/k)^{kf}$), using $g>0$ (as $g\ge(kb)^8>0$), $v\ge0$ and
$\Xi(s)\ge0$, we get $N_s(C,\Gamma)\le N_s(C^\sharp,1)$ and $G_s(C,\Gamma)\le G_s(C^\sharp,1)$ for $2\le s\le k$. Summing and applying the two
bounds for $C^\sharp$ gives the claim.
\end{proof}

\begin{proposition}\label{pf:open0:coll}
$\mathrm{On}_c(\mathrm{CO})$ holds for every $c\in\R$.
\formallabel{GlobalTheory.Core.Deg.collS_on}\ \formallabel{GlobalTheory.T.collisionAbsSum_bound_of_two_rowBounds}\
\formallabel{GlobalTheory.Core.Deg.collS_of_rowBounds}
\end{proposition}

\begin{proof}
\emph{Step 1: the class bounds.} Let $C\ge408\,e\pi$ and $C_T\le C$. For $2\le s\le k$, the non-degenerate and degenerate row
functions $N_s$, $G_s$ of \cref{pf:open0:corank} are evaluated at $(C,k,2r_S+2,f,s,b,g,v,\Gamma_S,\Xi)$, where $g=g_0g_1$, $v$ is
half the root discriminant of the host, $\Gamma_S\in[0,e^{k/8}]$ is a factor of the size data and $\Xi$ an explicit size function of
the words. The explicit size hypothesis (\cref{pf:open0:xisizes}, applied to the pair $(u,u)$, with the Siegel mass ratio) and
\cref{pf:open0:corank} control the sums over $s$.

\emph{Step 2: collisions from row bounds.} Let $K_0$ be the constant of \cref{pf:open0:corank} for $C$. Let $(X,Y)$ be concretely
tied and $u$ a word, with $\mathrm{H}_0(C_0)$, $C'\ge\max(K_0,4)$,
$C_l\ge\max(K_0,4)$, absolute convergence, a term tie $T$ and a loop tie $Z$ at $u$ with the tie conditions, the form condition at $u$,
and the non-degenerate and degenerate row bounds $\mathrm{R}_{\mathrm{nd}}(T,Z,C)$, $\mathrm{R}_{\mathrm{deg}}(T,Z,C)$. The count hypotheses for the
diagonal pair hold (\lref{Formalism.MainNotions.count_hypotheses_diagonal_words}); Step~1 gives the size hypothesis
$\mathrm{XH}(C)$, the loop-cut lemma gives the parameter conditions $\mathrm{PC}(C,K_0)$ with $m=2r_S+2$ and $N_T=T_u^f$
(\lref{Formalism.L2.SG.LoopsCut.params_at}), and $\Gamma_S\le e^{k/8}$ follows from
$8C_\Gamma\le C'$ and $C'm^5fL\le k$. For each $s$ with $2\le s\le k$, the term and loop ties bound the amplitude sum of the class
$\openzerobig_s$ by the rows of the loop tie (\lref{Formalism.L2.SG.LoopTieOf.classS_le}):
\[
\sum_{(t,t')\in\openzerobig_s}\|\mathrm{vec}\,t\|\,\|\mathrm{vec}\,t'\|\le M\,\pi_u\,\bigl(N_s\,r+G_s\bigr),\qquad r=d_u/D_*,
\]
where $M$ is the mass of the row representatives of the loop tie and $\pi_u$ a dimension factor of $u$. The constant $\Gamma_1$ of
$\mathcal T(u,u)$ equals a constant $\gamma_G$ of $D$ (\lref{Formalism.L2.SG.LoopTieOf.gammaFixed_of_hypotheses}), and the form
condition at $u$ states $\gamma_G\,M\,\pi_u=d_u$; hence (\lref{Formalism.L2.SG.LoopTieOf.gam_mul_le}) the hypothesis of
\cref{pf:open0:classes} holds with $B(s)=d_u(N_sr+G_s)$. Coranks are at most $k$. By \cref{pf:open0:classes} and
\cref{pf:open0:corank},
\[
\openzeropa(\openzerobig(\openzeroNZ))\le d_ur\sum_sN_s+d_u\sum_sG_s\le(g_0g_1)^{-kf/2}\,dd(u)+e^{-\openzerokfk/2}\,d_u,
\]
since $d_ur=dd(u)$. This is (CO).

\emph{Step 3: the hypotheses of Step~2.} Fix the width $c$ and the constants. The following four inputs hold uniformly: amplitude
bounds at the size data (\cref{pf:open0:ampon}),
bounds for the mean number of flag pairs of corank two (\lref{GlobalTheory.Core.Deg.corankTwoCountBounds_of_rowCountBoundsUniform},
from the Siegel mass ratio of \cref{prop:families:mass}, a theorem of the development
(\lref{GlobalTheory.Mass.siegelMassRatio_holds}), via \lref{GlobalTheory.Core.Deg.cntOn_singleton}),
smallness conditions (\lref{GlobalTheory.Core.Deg.smallOn_univ_unconditional}, restricted to the width by
\lref{GlobalTheory.Core.Deg.inputs_of_subset_widths}), and term and loop ties with the volume bound
(\cref{pf:open0:tieon}). From them one chooses $c^\sharp=\max(c_3,1,C_T,|c_H|+2,|c_V|+2)$ and $C$ as
the maximum of $c^\sharp$, $4+((c^\sharp)^2+2\max(C_h-c^\sharp,0))(6c^\sharp+2)$, $2C_2$, $|C_a|$, $408e\pi$ and $C_T$ (where $c_3,c_H$, $c_V$
and $C_a,C_2,C_h$ are the constants of the count, tie and amplitude inputs); the degenerate and non-degenerate row bounds at the
size data then hold (\lref{GlobalTheory.T.degenerateCollision_rowBound_at_sizeData} and \lref{GlobalTheory.T.nondegenerateCollision_rowBound_at_parameters}),
and Step~2 applies for $C',C_l$ beyond the maximum of the thresholds of the inputs, of $K_0$, of $3|C_2|$ and of $4$.
\end{proof}

\paragraph{The corank-one bounds.}
The two corank-one rows of the development are the following uniform statements for pairs of words.
\begin{itemize}
\item The \emph{generic row} $\mathrm{R}^{\mathrm{gen}}$: for all $C_g,c_b,C_w,c_T,C_T,C_0$ with $c_b,c_T>0$ there are $C_B\ge0$ and
thresholds such that, for all $C',C_l$ beyond them, all $D,L_2$, concretely tied $(X,Y)$ and words $u$ with the ties of $D$,
$\mathrm{H}_0(C_0)$, the standard conditions for $(u,u)$ and absolute convergence, (G1)$_{\gamma_L(C_B)}$ holds.
\item The \emph{non-generic row} $\mathrm{R}^{\mathrm{ng}}$ (\cref{prop:families:volratio}(iii)): the same with words $u,u'$,
a constant $C_4\ge0$, and the conclusion $\openzerops(\openzeroone(\openzeroNG))\le\varepsilon_{\mathrm{NG}}(C_4)\,dd(u,u')$ for
$\mathcal T(u,u')$, where $\varepsilon_{\mathrm{NG}}(C_4)=2^{-(k-2r_S-5)}\exp(C_4m^5f\log b)$.
\end{itemize}

\begin{lemma}\label{pf:open0:onerows}
$\mathrm{R}^{\mathrm{gen}}$ implies $\mathrm{One}^{\mathrm{gen}}_c$, and $\mathrm{R}^{\mathrm{ng}}$ implies $\mathrm{One}^{\mathrm{ng}}_c$, for every $c\in\R$.
\formallabel{GlobalTheory.Core.Diag.oneOn_singleton_of_rows}\ \formallabel{GlobalTheory.Core.Diag.epsNG_le_gamL}
\end{lemma}

\begin{proof}
Given the constants, apply the row with $C_0=c$; it gives a constant ($C_B$, resp.\ $C_4$) $\ge0$ and thresholds $C'_0,C_{l,0}$. Take this
constant for $C$ and $t=\max(C'_0,C_{l,0},3)$. For $C',C_l\ge t$ the hypotheses of $\mathrm{One}_c$ give those of the row: since $C_l>2$,
$\mathrm{H}_0(c)$ implies the standard conditions for the pair $(u,u)$ (\lref{Formalism.MainNotions.Stand_of}). For the generic row
the conclusion is (G1)$_{\gamma_L(C_B)}$ (the constant of the development equals $\gamma_L(C_B)$). For the non-generic row,
$\log b=L$ and $2^{-(k-2r_S-5)}\le1$ give $\varepsilon_{\mathrm{NG}}(C_4)\le\gamma_L(C_4)$, and $dd(u)\ge0$, so
(N1)$_{\gamma_L(C_4)}$ holds.
\end{proof}

\begin{proof}[Proof of \cref{pf:open0:fourpoints}]
(a) The degenerate bound is \cref{pf:open0:deg} and the collision bound is \cref{pf:open0:coll}; the Siegel mass ratio used in the latter
is a theorem of the development (\lref{GlobalTheory.Mass.siegelMassRatio_holds}).

(b) Assume the Siegel--Weil formula. Then $\mathrm{R}^{\mathrm{ng}}$ holds (\lref{GlobalTheory.rowNGP_of_siegelWeil}). For
$\mathrm{R}^{\mathrm{gen}}$: the Siegel--Weil formula gives the volume identities $\sfname{SWXS}$ (\lref{GlobalTheory.St.swxs_of_siegelWeil})
and, together with the real objects at infinity (\lref{Formalism.L2.SG.RealObj.EnergyConvJ.BM.realObj_inInfinityX_jsz}, using the
uniqueness of Bessel models with large gap \lref{CompanionTheory.Lime.BM.uniqGap_holds}), the step choice constants
(\lref{GlobalTheory.St.stepChoiceDataC_of_siegelWeil_inInfinity}). With the theorem-B step for the mass data
(\lref{Formalism.MainStatements.RowRest.bDat_theoremB_global}), the estimates for equal words of the non-generic chain
(\lref{GlobalTheory.T.equalWords_nonGeneric_estimates_chain}: the point, norm-one, tie and pin steps) and the normal count,
line, integrality and separation statements (\lref{GlobalTheory.T.normalCount_line_integrality_separation}), the generic
corank-one row theorem (\cref{pf:open0:gen1}) gives $\mathrm{R}^{\mathrm{gen}}$. \Cref{pf:open0:onerows} concludes.
\end{proof}

\begin{remark}\label{pf:open0:rem:diag}
(a) Of the Siegel--Weil formula, the proof of \cref{pf:open0:diagthm} uses only the corank-one bounds; the degenerate and collision
bounds hold without it. (b) The truth of $D$ in the strong sense is used only to obtain spectral ties with finite spaces and local
ties; the identity condition of the Hida setting of the world, which distinguishes the strong sense, is not used here. (c) The
constant $C$ of \eqref{pf:open0:eq:diag} is $\max(C_1,C_2)+1$, where $C_1,C_2$ come from the two corank-one rows; the factor $4$ comes
from $\frac32\cdot2\cdot\frac65=\frac{18}5\le4$. (d) The width $3$ is the constant $C_0=3$ of the dimension window of
\cref{pf:open0:window}(c), which follows from the rank identity and the window for the level prime (in fact with the better
interval $[\frac12,3]$); the four point bounds are used only on this width.
\end{remark}

%==============================================================================
\subsection{The definite-side row}\label{sec:pf:open0:df}

The definite-side condition is a height bound for the column chains of the words of a bundle, under the point inequalities
(\cref{sec:families:nonempty}). Its formal statement quantifies over the data produced by the point inequalities; we state
it in this form.

\begin{definition}[Definite-side condition]\label{pf:open0:def:df}
Let $b$ be a bundle over $(D,O,c)$ and $C,C_H\in\R$. $\openzeroDF(D,O,c,b,C,C_H)$ is the statement: for all proofs of the
first-level admissibility $\mathcal O_3(1)$ and of the separation of types of $b.P$, all covolume-onto-isometry conditions and all
pass-in conditions for the realisation of $b$ with constants $(C,C_H)$, all reals $\lambda_0,C_0,c_e,S_w,C_1$ with the
numerical facts for the recursive numbers of every word, all families of $\OO_{\mathit{LL}}$-lattices $N^C_w$, all
three-sided division compatibility data $I$ which are tied to the package built from the primitive links and equal to the
division data of $b.P$, and every word $u$, the \emph{column rows} hold for $u$:
\begin{itemize}
\item there are a column chain $\mathrm{ch}$ of the face for the member $u$, with a chain tie and its first columns;
\item a frame $X_0$ at the infinite places $v$ of $\mathit{LL}$ with $(X_0(v)^{-1})^*X_0(v)^{-1}$ equal to the image under $v$ of the
hermitian metric of the base row space at $\mu_0$;
\item height numbers $h$ for the recursive numbers of $u$ such that, at every $v$, the square of the operator norm of the final
point map of $\mathrm{ch}$ (from the standard basis to the dual frame of $X_0$) is at most $e^{h_{\mathrm{Pn}}}$;
\item Cauchy--Binet data for the chain whose squares satisfy $\log\mathrm{cb}(j,v)\le\sum_{l\mapsto j}h_{\mathrm{sz}}(l,v)$;
\end{itemize}
moreover the column metric of the member $u$ is invertible at every $v$, the number $a(u)$ of columns is positive, and the
lattice $N_u$, transported to the column space, is a lattice.
\formal{GlobalTheory.Fc.DefiniteSideCondition}
\end{definition}

\begin{definition}[Point-evaluation side]\label{pf:open0:def:pes}
A bundle $b$ over $(D,O,c)$ has the \emph{point-evaluation side} if: $D$ is real-tied and admits symmetric tied column forms whose
step-zero characters are finite; the column characters of $b$ are compatible; the column weights of $b$ form, after the
identification of index types, a system of weights $W$ for the kernel datum of the record $b.R$ such that the kernel ties
$\mathrm{KT}(t,W,\Delta,h_0)$ of \cref{pf:open0:def:kt} hold for its kernel entry $t$; and for all $C_Y,C_T$, if the height rows of the size
data hold with $C_Y$, $X'\ge8$, $x\ge2$, $4a<T_w$ and $\log T_w\le C_TmL$ for every word $w$, then the kernel sizes
$\mathrm{KS}(t,W,\Delta,h_0,C_Y,C_T)$ hold.
\formal{GlobalTheory.Fc.RL.DefSide.PointEvaluationSide}
\end{definition}

\begin{proposition}\label{pf:open0:pes}
For every parameter $p$, every bundle linked to a point by $\openzeroLk_p$ has the point-evaluation side.
\formallabel{GlobalTheory.Fc.RL.DefSide.pointEvaluationSide_of_isLinked}\
\formallabel{GlobalTheory.Fc.RL.DefSide.pointEvaluationSide_of_bWit}
\end{proposition}

\begin{proof}
It suffices to treat a bundle $b$ with a witness $w$ over a weighted point $q$. The point is real-tied, and the member
construction provides symmetric tied column forms with finite step-zero characters
(\lref{GlobalTheory.T.exists_symmetricColumnForms_stepZeroFiniteCharacter}). The column pins of the witness are the
compatibility of the column characters. The kernel tie of the columns of $b$ identifies the index types of its weights with
those of the kernel datum of $b.R$, and the witness identifies the weights with those of the point. The kernel datum of $b.R$
is standard at the point (a component of the record compatibility of the witness), so \cref{pf:open0:kerties} gives the kernel
ties, and \cref{pf:open0:kersizes} gives the last condition.
\end{proof}

\begin{proposition}\label{pf:open0:dfpoint}
For all $C_Y\in\R$ and $C_T\ge0$ there is $C_0$ with the following property. Let $b$ be a bundle over $(D,O,c)$ with the
point-evaluation side, such that the height rows of the size data hold with $C_Y$, $X'\ge8$, $x\ge2$, $4a<T_w$ and $\log T_w\le C_TmL$
for every word $w$, the condition $N'$ of the denominator datum holds, the lattices $N_w$ have full rank and positive rank, and
$c\ge C_0$. Then $\openzeroDF(D,O,c,b,C,C_H)$ for all $C,C_H$.
\formallabel{GlobalTheory.Fc.RL.DefSide.definiteSideCondition_of_nprimeDef_of_nuFull}
\end{proposition}

\begin{proof}
Let $C_0$ be the constant of the kernel column theorem (\cref{pf:open0:kercols}) for $(C_Y,C_T)$. Given the data quantified in
$\openzeroDF$, the package built from the
primitive links of $b$ has the same $Z_s$, base and columns as $b.P$, and the tie of $I$ with this package and $I=b.P.I$ identify the
member of $I$ at $u$ with $u$. The point-evaluation side gives the remaining hypotheses, and the kernel column theorem with
$C_h=c\ge C_0$ gives the column rows.
\end{proof}

\begin{proposition}[The definite-side row]\label{pf:open0:df}
For every parameter $p$, $\openzeroLinked(p,\mathsf{df})$. (The formal statement carries four classical results as
hypotheses: Weil representations, existence of discrete series, finiteness and rapid decay; they are not used.)
\formallabel{GlobalTheory.T.weil_unitary_entries'}\ \formallabel{GlobalTheory.Fc.RL.linkedAt_prmR_df_of_cuspRapidDecay}\
\formallabel{GlobalTheory.Fc.RL.DefSide.linkedAt_df_of_parameters}
\end{proposition}

\begin{proof}
By \ref{pf:open0:R4} it suffices to prove $\openzeroRow(\openzeroGood_p,\openzeroLk_p,\openzeroEvWC(B^{\mathrm{df}}))$. Take for $E$ the
point-evaluation side; $\openzeroEv(E)$ holds by \ref{pf:open0:R2} and \cref{pf:open0:pes}. Let $\sigma$ be an $E$-selection over a
class $\mathcal A$. For every $(C_Y,C_T)$ let $c_0(C_Y,C_T)$ be the constant of \cref{pf:open0:dfpoint} for $(C_Y,\max(C_T,0))$.
Closedly on every class, $c_0(C_{YL},C_{Tu})\le C_{\mathrm{sz}}$, where $C_{YL}$ and $C_{Tu}$ are the entries $6$ and $13$ of the
constant list of the family and $C_{\mathrm{sz}}$ its size constant (\lref{Formalism.MainNotions.DIV.closed_Csz_ge_fn}); we
prove $\openzeroEvWC(B^{\mathrm{df}})$ closedly under this condition. Let $\mathcal F$ be a family of $\mathcal W_3(\mathcal A)$ with
$c_0(C_{YL},C_{Tu})\le C_{\mathrm{sz}}$. The selection provides a row parameter linked to $\mathcal F$ and, for $x>x_1$, that
$\sigma(\mathcal F,x,i)$ is a built bundle over $(D_{x,i},O_{x,i},C_{\mathrm{sz}})$ with the point-evaluation side. Its first component
gives the real-place second-layer conditions of the family (for $x>x_1+1$); with the size condition $\sfname{TuSize}$ and the
chosen-$E$ data of the family (both part of membership in $\mathcal W_3(\mathcal A)$), \cref{pf:open0:fullrank} gives thresholds $a_1,a_2$ beyond which the lattices $N_w$ have full rank and positive
rank, and \lref{GlobalTheory.CoreHt.famRowsSizes} gives a threshold $a_3$ beyond which the height rows with $C_{YL}$, $X'\ge8$,
$x\ge2$, $4a<T_w$ and $\log T_w\le\max(C_{Tu},0)\,mL$ hold whenever $\Delta$ is a natural number (which the bundles guarantee). The
condition $N'$ is a field of the family. Take $C_0=0$ in the definition of $\openzeroEvWC$ and $x_0=\max(a_1,a_2,a_3)$: for $C\ge0$,
$x\ge x_0$ and $i\in I_x$, \cref{pf:open0:dfpoint} applies with $c=C_{\mathrm{sz}}\ge c_0$ and gives $B^{\mathrm{df}}_{C,C_H}$ at
$\sigma(\mathcal F,x,i)$.
\end{proof}

%==============================================================================
\subsection{Proofs of the inputs}\label{sec:pf:open0:inputs}

This subsection proves the inputs used above whose formal proofs belong to the part of the development converted in this
section: the ranks and the kernel column theorem behind the definite side (\cref{sec:pf:open0:ranks,sec:pf:open0:kernel}),
the type labels, the label ties, the non-vanishing of isotypic components and the identification of $\kappa$ behind the label
and $\kappa$-rows (\cref{sec:pf:open0:labels,sec:pf:open0:kappa}), the division compatibility (\cref{sec:pf:open0:divc}), and
the identity cells, the pair amplitude bound, the inputs of the collision bound and the generic corank-one row of the
diagonal row theorem (\cref{sec:pf:open0:diaginputs}). Each proof assembles deeper results of the development, which are
stated and cited.

\subsubsection{Ranks of the definite side}\label{sec:pf:open0:ranks}

Let $\mathcal F$ be a family (\cref{def:families:closed}) with members $D_{x,i}$ for $x>x_1$. For a property $P$ of pairs
$(D,u)$ of a datum and a word of it, we write $\mathrm{Fam}_{\mathcal F}(P)$ if there is $x_0$ such that $P(D_{x,i},u)$ holds for
all $x\ge x_0$ with $x>x_1$, all members $i$ and all words $u$ of $D_{x,i}$ carrying a Hida setting. Recall the properties
(G7)--(G11) of \cref{pf:top-shape:famglmin}: free $T$-data, the slope line, Hecke integrality at level zero, ramified level
zero and the second-branch inclusion length. For a word $u$ let $d_U(u)$ and $d^S_u$ be the ranks of its member and of its
$S$-part, $M_u(0)$ its level-zero lattice and $M_u^{\mathrm{ram}}$ its ramified lattice, with mean slopes $\mu$ and maximal
slopes $\mu_{\max}$. Consider the two properties:
\begin{itemize}
\item $P_2(D,u)$ (\emph{property two}): $d_U(u)=d^S_u\bigl(1-\frac1{T_u-1}\bigr)^f$,
$\mu(M_u^{\mathrm{ram}})=\mu(M_u(0))-\frac{f\log(T_u-1)}{T_u-2}$ and $\mu_{\max}(M_u^{\mathrm{ram}})\le\mu_{\max}(M_u(0))$;
\item $P_5(D,u)$ (\emph{property five}, the congruence-lattice length bound): at every place $w$ above $T_u$ the lattice
$N_\lambda(w)$ is contained in the ramified lattice $S_{\mathrm{ram}}(w)$, and the length of $S_{\mathrm{ram}}(w)/N_\lambda(w)$ is at
most $e_w\,m\,d^S_u\sum_v\bigl(a_{\max}/k+2\,\mathrm{wt}(\tau_v)+k+b_{\mathrm{top}}(\tau_v)\bigr)$, where $e_w$ is the ramification
index of $w$ and $v$ runs over the column places above $w$, with types $\tau_v$, weights $\mathrm{wt}$ and top bounds
$b_{\mathrm{top}}$.
\end{itemize}

\begin{proposition}[Properties two and five]\label{pf:open0:twofive}
Let $\mathcal F$ be a family. If $\mathrm{Fam}_{\mathcal F}$ holds for (G7) and for (G8), then $\mathrm{Fam}_{\mathcal F}(P_2)$; if
$\mathrm{Fam}_{\mathcal F}$ holds for (G9), (G10) and (G11), then $\mathrm{Fam}_{\mathcal F}(P_5)$.
\formallabel{Formalism.MainNotions.Wire.family_theorems_two_and_five}
\end{proposition}

\begin{proof}
If $\mathrm{Fam}_{\mathcal F}(P)$ and $\mathrm{Fam}_{\mathcal F}(Q)$ hold with thresholds $x_0$ and $x_0'$, then
$\mathrm{Fam}_{\mathcal F}(P\wedge Q)$ holds with threshold $\max(x_0,x_0')$; and if $P$ implies $Q$ at every word carrying a Hida
setting, then $\mathrm{Fam}_{\mathcal F}(P)$ implies $\mathrm{Fam}_{\mathcal F}(Q)$. It therefore suffices to prove the two
implications at a single word $u$ of a datum $D$ carrying a Hida setting. For $P_2$: a free $T$-datum for $u$ (G7) gives the
rank identity and the bound for $\mu_{\max}$, and the slope line (G8) is the identity for $\mu$
(\lref{Formalism.MainNotions.product_first_zeroth_parts.division_side_conditions_at_sizeData}). For $P_5$: at a word with a
Hida setting, the second-branch inclusion length, Hecke integrality at level zero and the ramified level-zero condition imply
the congruence-lattice length bound (\lref{Formalism.L2.GlobalTheoryTheory.LocalLatticeBases.congruence_no_kernel_minimal}).
\end{proof}

\begin{lemma}[Ranks]\label{pf:open0:fullrank}
Let $\mathcal F$ be a family which satisfies the real-place second-layer conditions eventually, $\sfname{TuSize}$ and
$\sfname{ChosenE}(3,C_H)$ for some $C_H$. Then eventually the lattices $N_u$ of all words of all members of $\mathcal F$ have
full rank ($\sfname{NuFullF}$) and positive rank ($\sfname{MuRank}$).
\formallabel{GlobalTheory.Fc.RL.DefSide.full_rank}\ \formallabel{Formalism.MainNotions.muRank_items}
\end{lemma}

\begin{proof}
The second-layer conditions give, in the form $\mathrm{Fam}_{\mathcal F}$, the concrete forms at level zero (the twist
compatibility of the forms together with an admissibility condition) (\lref{Formalism.MainNotions.family_conditions_of_layer_two}).
Free $T$-data, the slope line and the second-branch integrality under forms hold at every word of every datum
(\lref{Formalism.MainNotions.freeTT_holds}, \lref{Formalism.MainNotions.slopeLineT_holds}, and \lref{Formalism.L2.GlobalTheoryTheory.LemB.second_branch_all_members});
together with the concrete forms they give (G7), (G8) and (G11) in the form $\mathrm{Fam}_{\mathcal F}$
(\lref{Formalism.MainNotions.RT.slope_free_branch_family}). The integrality at the base point lattice, which holds at every
word (\lref{PI.IntegralAtBasePointLattice.integral_at_t0_lattice}), and the concrete forms give (G9)
(\lref{Formalism.MainNotions.RT.hI_tied}), and the concrete forms at level zero give (G10) at every word with a Hida
setting (\lref{Formalism.L2.GlobalTheoryTheory.UnramifiedTrace.hecke_ramified_zero_all}). By \cref{pf:open0:twofive},
$\mathrm{Fam}_{\mathcal F}(P_2)$ and $\mathrm{Fam}_{\mathcal F}(P_5)$ hold.

Now $\sfname{TuSize}$ gives the eventual level-prime setting (\lref{Formalism.L2.GlobalTheoryTheory.Places.settingF_of}); with
$P_2$ and $P_5$ it gives the level statement $\sfname{LvlF}$ (\lref{Formalism.MainNotions.lvlF_of}), which implies the
commensurability of the lattices $N_u$ and hence $\sfname{NuFullF}$
(\lref{Formalism.MainNotions.nuCommensF_min} and \lref{Formalism.MainNotions.Lnk.nuFullF_of}). Finally $\sfname{ChosenE}(3,C_H)$
contains the proviso $\sfname{ChosenProviso}(3)$; since $3>0$, the proviso, $\sfname{NuFullF}$, $\sfname{LvlF}$ and the setting give
$\sfname{MuRank}$ (\lref{Formalism.MainNotions.muRank_min}).
\end{proof}

\subsubsection{Kernel ties, kernel sizes and the kernel column theorem}\label{sec:pf:open0:kernel}

The point-evaluation side (\cref{pf:open0:def:pes}) refers to the kernel ties and the kernel sizes of a kernel entry. At a point
of the construction the columns come from a construction $I$ of \emph{true columns} over a base with weights $W_B$, together
with laws $l$ and later data $Z$; the record of the point is $D=D(I,l,Z)$. A \emph{tied kernel entry} $t$ over an extra datum of
$D$ consists of Hida fields $X_t$, a seam, primes, characters, a selection of places, global character data, words, tuples and
constants at every tuple, an ambient datum, analytic data $(\delta,\chi,c_D)$ and an integer $m_2$, subject to compatibility
conditions; it is \emph{standard} for $(I,l,Z)$ if these components are those of the true columns of $I$. For size data $L$ of
$D$ and $\Delta,h_0\in\N$ let $N=N(D,L,\Delta,h_0)$ be the numerical datum of the record, and for a word $u$ let $\delta_u$ be its
code in $N$.

\begin{definition}[Kernel ties and kernel sizes]\label{pf:open0:def:kt}
Let $t$ be a tied kernel entry, $W$ weights for its primes and $\Delta,h_0\in\N$. The \emph{kernel ties}
$\mathrm{KT}(t,W,\Delta,h_0)$ are the following conditions.
\begin{enumerate}[label=(T\arabic*)]
\item\label{pf:open0:T1} (Host.) The record, the invariants $q,d$, the type primes and the host axioms of $N$ are those of
the host of $D$, whose class datum is full, whose CM type is induced and which satisfies the host field conditions; the fields
$X_t$ satisfy the unramified abelian extension conditions; the seam and the primes are admissible, the base embedding
$\sigma_0$ commutes with complex conjugation, and the primes satisfy the $\pi_G$-congruence and have the conductor and the
generators of $N$.
\item\label{pf:open0:T2} (Weights.) $W$ is admissible, and there is a bijection $p$ from the real places of $N$ to the
embedding types of $X_t$ under which the top places correspond and, at every word $u$ and real place $z$, the row parameter and
the weights of $W$ at $(u,p(z))$ are those of $N$ at $(z,\delta_u)$.
\item\label{pf:open0:T3} (Characters.) The characters of $t$ are admissible, tied to the row weights of $W$ through a Hecke
character and Weil-tied; the base embedding and the character $\xi_0$ of the global data match the seam and the host; the
selection of places is a selection.
\item\label{pf:open0:T4} (Analytic.) $\delta$ is a twelfth root of the discriminant of the ideal lattice, $\chi$ is the
character of $D$ and $c_D$ is its $d$-generator.
\item\label{pf:open0:T5} The tuple of $t$ at the base tuple $\mu_0$ is trivial, the constant of $t$ at $\mu_0$ is sharp, and
$m_2=m-2$.
\item\label{pf:open0:T6} (Field of definition.) For every word $u$ and every $i<m$, the jets of the $i$-th chain step of $u$
at $\mu_0$ are defined over the coefficient field.
\end{enumerate}
For $C_Y,C_T\in\R$ the \emph{kernel sizes} $\mathrm{KS}(t,W,\Delta,h_0,C_Y,C_T)$ are: (S1) the size conditions of the primes and
weights with the parameter $a$ of $L$; (S2) for every word $u$ the \emph{size rows} $\mathrm{SR}_{C_Y}(N,\delta_u)$, a list of
explicit inequalities, among them: the logarithms of the root discriminant, of $q$, of $d$ and of $4x$ and the logarithmic
derivative at $1$ are at most $L$, $(kb)^8\le g_i$ and $\log g_i\le C_YL+1$, $f\,\mathrm{Nm}\le b$, and the $\nu$-weights, the
differences of Harish-Chandra parameters, the weights and the top bounds are bounded by $3a$, $\eta$, $S_c$ and $2b'$; (S3)
$\log T_u\le C_TmL$ for every word $u$.
\formal{GlobalTheory.ColRows.KerTies}\ \formal{GlobalTheory.ColRows.KerSizes}\ \formal{GlobalTheory.Core.Ht2.SizeRowsS}
\end{definition}

\begin{proposition}[Kernel ties]\label{pf:open0:kerties}
Let $I$ be an admissible construction of true columns with pins and with jets in the coefficient field, $l,Z$ laws and later
data, $D=D(I,l,Z)$ with size data $L$ and an extra datum, $t$ a tied kernel entry standard for $(I,l,Z)$, $W=W_B$, and
$\Delta,h_0\in\N$ such that $\Delta$ is the parameter $\Delta$ of $D$ and $mh_0$ the $\eta$-size of $L$. If the constant of $t$ at
$\mu_0$ is sharp, then $\mathrm{KT}(t,W,\Delta,h_0)$.
\formallabel{GlobalTheory.T.kernelCompatibilities_standardKernelData_trueColumns}\
\formallabel{GlobalTheory.ColRows.kerTies_of_stdKer}\
\formallabel{GlobalTheory.T.hostCompatibility_standardKernelData_trueColumns}\
\formallabel{GlobalTheory.ColRows.hostTie_stdKer}\
\formallabel{GlobalTheory.T.weightCompatibility_standardKernelData_trueColumns}\
\formallabel{GlobalTheory.ColRows.weightTie_of_stdKer}\
\formallabel{GlobalTheory.T.characterCompatibility_standardKernelData_trueColumns}\
\formallabel{GlobalTheory.ColRows.charTie_stdKer}\
\formallabel{GlobalTheory.T.analyticCompatibility_standardKernelData_trueColumns}\
\formallabel{GlobalTheory.ColRows.analyticTie_of_stdKer}\
\formallabel{GlobalTheory.T.coefficientField_isFieldOfDefinition_chainSteps}\
\formallabel{GlobalTheory.ColRows.fieldOfDefinition_of_stdKer}
\end{proposition}

\begin{proof}
By standardness every component of $t$ is the corresponding component of the true columns of $I$, so each condition is a
statement about the true columns. \ref{pf:open0:T1} is the host tie of the true columns
(\lref{GlobalTheory.CoreHt.host_tie_trueCols}). For \ref{pf:open0:T2}: since $\Delta$ is the parameter $\Delta$ of $D$ and
$mh_0$ the $\eta$-size, the numerical datum attached to the size data equals $N(D,L,\Delta,h_0)$
(\lref{GlobalTheory.T.columnNormParameters_eq_natStepGap}), so \ref{pf:open0:T2} is the weight tie of the true columns for
the bijection of real places of the record (\lref{GlobalTheory.CoreHt.weightTie_trueCols}). \ref{pf:open0:T3} is the
character tie of the true columns (\lref{GlobalTheory.CoreHt.charTie_trueCols}). In \ref{pf:open0:T4} the first two
conditions are part of the tied kernel entry; the $d$-generator $c_D$ of $t$ and the $d$-generator of the true columns are both
the $d$-generator of the objects at the base word, and the latter has the required property by the analytic ties of the true
columns (\lref{GlobalTheory.Core.Ht2.analyticTies_trueCols}). In \ref{pf:open0:T5} the tuple of the true columns at $\mu_0$ is
trivial, sharpness is assumed and $m_2=m-2$ is part of standardness. For \ref{pf:open0:T6}, the analytic ties of the true columns
also give the field of definition for the character $\chi_I$ of the true columns; $\chi_I$ and $\chi$ agree on the determinants of
all standard lattices (\lref{GlobalTheory.ColRows.FldDef.chiRec_detId}), and the field-of-definition condition depends on the
character only through these values (\lref{GlobalTheory.ColRows.FldDef.fieldOfDefinition_congr}).
\end{proof}

\begin{proposition}[Kernel sizes]\label{pf:open0:kersizes}
In the situation of \cref{pf:open0:kerties} (without the hypotheses on jets and sharpness), let $C_Y,C_T\in\R$ and assume the
height rows of $L$ with $C_Y$, $X'\ge8$, $x\ge2$, $4a<T_u$ and $\log T_u\le C_TmL$ for every word $u$. Then
$\mathrm{KS}(t,W,\Delta,h_0,C_Y,C_T)$.
\formallabel{GlobalTheory.ColRows.kerSizes_of_stdKer}
\end{proposition}

\begin{proof}
(S1) holds for the true columns under the height rows, $X'\ge8$, $x\ge2$ and $4a<T_u$
(\lref{GlobalTheory.T.trueColumns_prime_weight_sizes}), and the primes and weights of $t$ are those of the true columns by
standardness. (S2): the height rows imply the size rows at the numerical datum attached to the size data for every word
(\lref{GlobalTheory.T.sizeInequalitiesNoLevelPrime_columnNormParameters}), and this datum is $N(D,L,\Delta,h_0)$ as in the proof
of \cref{pf:open0:kerties}. (S3) is assumed.
\end{proof}

\begin{theorem}[Kernel column theorem]\label{pf:open0:kercols}
For all $C_Y\in\R$ and $C_T\ge0$ there is $C_0$ with the following property. Let $D$ be a datum with size data $L$ and an extra
datum whose field $\mathit{LL}$ is totally complex, $T$ tied data with kernel entry $t$, $P$ a row-side family and $K$
square-integrable columns for $(T,P)$ with compatible column characters and column weights $W$. Let $\Delta,h_0\in\N$ with
$\mathrm{KT}(t,W,\Delta,h_0)$ and $\mathrm{KS}(t,W,\Delta,h_0,C_Y,C_T)$, and assume the condition $N'$ of the denominator datum and
that all lattices $N_w$ have full rank and positive rank. Let $\mathrm{Pk}$ be a scaled package of the face of $K$ with the same
$Z_s$ as $P$, $u$ a word with the numerical facts for $(\lambda_0,c_0,c_e,C_h,S_w,C_1)$ at its code, and
$e\colon\mathit{LL}^{a(u)}\to V(N_u)$ an $\mathit{LL}$-linear isomorphism onto the space of $N_u$. If $C_h\ge C_0$, the column rows of
\cref{pf:open0:def:df} hold for $u$.
\formallabel{GlobalTheory.T.kernel_columns_from_ties_alt}\ \formallabel{GlobalTheory.ColRows.rowsU_of_colPins_alt}\
\formallabel{GlobalTheory.T.memberRows_of_kernelDatum_sizes}\ \formallabel{GlobalTheory.ColRows.rowsBody_of_kerTies}\
\formallabel{GlobalTheory.T.columnsMemberRows_of_logLevel_le}\ \formallabel{GlobalTheory.ColRowsTopS.rowsBody_colFace}\
\formallabel{GlobalTheory.ColRows.exists_isFrames}
\end{theorem}

\begin{proof}
\emph{Step 0: the metric, the lattice and the rank.} The hermitian metric $H=H_V(\mu_0)$ of the base row space is positive
definite at every infinite place $v$ of $\mathit{LL}$ (\lref{Formalism.L2.DIV.BlockDat.HV_pos}: a consequence of the
conjugation, level and irreducibility conditions of the columns); in particular it is invertible. The lattice $N_u$, transported
to the column space by $e$, is a lattice (\lref{GlobalTheory.ColRows.nuLat_comap_isLattice}, from $N'$ and the full rank of
$N_u$). Since $N_u$ has full rank, its rank is $\dim V(N_u)=a(u)$, which is therefore positive. Finally the frame values at $\mu_0$ are
bounded (\lref{GlobalTheory.T.frames_tilde_gamma_values_alt}: when the tuple at $\mu_0$ is trivial, for every frame $X_0$ of $H$,
every place $v$ and every column $F$, the evaluation of the inner product of the columns at $F$ is at most the square of the
$v$-norm of $F$ in the frame); the tuple at $\mu_0$ is trivial by \ref{pf:open0:T5}.

\emph{Step 1: identifications.} The kernel tie of the columns identifies the Hida fields, the record, the primes, the characters,
the selection, the support, the tuples, the constants, the indices, $m_2$, the global data, the ambient datum, the members, the
words and the level of $K$ with those of $t$, and the weights of $K$ with $W$; after these identifications the member of the face
at $u$ is $u$, and the hypotheses below are the components of $\mathrm{KT}$, $\mathrm{KS}$ and of the columns.

\emph{Step 2: the constant.} Let $C$ be the constant of the kernel-half bound for $C_Y$
(\lref{GlobalTheory.T.kernelHalf_chainNumbers_any_levelPrime}: there is $C\ge0$ such that, under the kernel ties, the size rows
and the root discriminant bound, the recursive numbers $R$ of the column chain of every word satisfy the kernel-side bound
$\mathrm{KS}_C(R)$, which bounds the logarithmic sizes of the lifted maps of the kernel half of the chain by
$Cfm^2N^2(\log g+\log T)+CfN^2\log g$ and similar expressions), and put
\[
C_0=\max\bigl(2(16+C_T),\ 6C+1+2C^2,\ 4C(C_T+C)\bigr).
\]

\emph{Step 3: height numbers.} Under the weight ties, the size rows, $\log T_u\le C_TmL$, $T_u\ge1$ and the root discriminant bound,
the recursive numbers of $u$ satisfy the last-side bound $\mathrm{LS}_{C_T}(R,A)$ for a real $A$
(\lref{GlobalTheory.T.lastMapHalf_recNums_of_log_levelPrime_le}): $A\le-\sum_\tau n_\tau\log T+\sum_\tau(2\log c_\tau+\log\xi_\tau)$
with explicit bounds for $n_\tau$, $c_\tau$ and $\xi_\tau$ in terms of $N$, $S$, $\eta$, $a$ and $b'$. Since $C_h\ge C_0$, the two bounds
combine into height numbers $h$ for $R$ with $h_{\mathrm{Pn}}=A$ and the size bounds required of Cauchy--Binet data
(\lref{GlobalTheory.ColRows.exists_htNums_of_kerSide_lastSide}).

\emph{Step 4: the frame.} For every infinite place $v$, the matrix $H_v$ is positive definite, so $H_v=A_v^*A_v$ with
$A_v=H_v^{1/2}$ invertible; $X_0(v)=A_v^{-1}$ satisfies $(X_0(v)^{-1})^*X_0(v)^{-1}=H_v$. This is the frame of the column rows.

\emph{Step 5: the chain.} The first columns of $u$ form an isometric embedding with non-zero determinant
(\lref{GlobalTheory.ColFirstIso.colsIn_isometry}); the weights of $u$ interlace the true weights, with the same shift
(\lref{GlobalTheory.BranchMembers.interlace} and \lref{GlobalTheory.IndefiniteRowSideFullHecke.InterlacingEquivariance.members_sameShift});
the evaluation data at $u$ are framed (\lref{GlobalTheory.ColRows.evGSD_datFramed}); the steps of the chain are framed
(\lref{GlobalTheory.FamilyFrames.famFramed_chF}); and the column periods are the periods of the column face
(\lref{Formalism.L2.DIV.colPer_eq}). Under these conditions the chain of column matrices of the column face at $u$, with the
frame $X_0$, the height numbers $h$ and their size bounds, satisfies the column rows, provided the final point map satisfies the
period bound of Step~6 (\lref{GlobalTheory.T.kerChain_matrices_with_norm_bounds}: the factor bounds of the steps and the
Cauchy--Binet bounds follow from the frames and the height numbers).

\emph{Step 6: the period bound.} For every place $v$, every $\psi$ and every column $F$, the square of the $v$-norm of the period of
the last map is at most $e^{A}$ times the evaluation of the inner product at $F$
(\lref{GlobalTheory.ColRowsTop.period_le}, from the norm of the last map, the interlacing and the same shift), and the latter is
at most the square of the frame norm of $F$ by Step~0. Since $h_{\mathrm{Pn}}=A$, this is the bound
$\le e^{h_{\mathrm{Pn}}}$ of the column rows. The invertibility of the metric, $a(u)>0$ and the lattice property are Step~0.
\end{proof}

\subsubsection{Type labels and label ties}\label{sec:pf:open0:labels}

For a datum $D$, the finite types $\sigma$ of $D$ carry labels $\mathrm{lab}_k(\sigma,v)$ at the places $v$ of $E$ above the type
primes $g_k$ ($k=0,1$). For a label $\lambda$ at $g_k$ let $e_{\mathrm{ser}}(k,\lambda)$ be the series idempotent attached to it. A
\emph{type representation} $\mathrm{rep}$ assigns to every finite type $\sigma$ a finite type $\mathrm{rep}(\sigma)$ of $\GL_r$ over $E$.
The \emph{type labels} $\mathrm{TL}(D,\mathrm{rep})$ are the conditions: for all $\sigma$, $k$ and $v$, the label
$\mathrm{lab}_k(\sigma,v)$ is a label, and at every split place $s$ above $v$ the finite type $\mathrm{rep}(\sigma)$ is fixed by
$e_{\mathrm{ser}}(k,\mathrm{lab}_k(\sigma,v))$.

\begin{theorem}[Type labels]\label{pf:open0:typelab}
Assume Green's theorem and Zelevinsky's theorem, and let $h_L$ be a proof of $\sfname{RecTExists}$. Let $a$ be the norm data of a
point over a datum $D$, with type representation $\mathrm{rep}$, such that the label tie of $D$ with the label data at the type
primes holds and the generator pin $\zeta_k=\mathrm{gen}_k$ ($k=0,1$) holds. Then $\mathrm{TL}(D,\mathrm{rep})$.
\formallabel{GlobalTheory.Fc.LabLinked.typeLab_of_gen_alt_alt}\ \formallabel{UpperTheory.Lit.Dz.localBesselMultiplicityOne}\
\formallabel{DiscreteGroups.sglob.DzType.Rt.rowTruth_of_labelTie_alt}\
\formallabel{DiscreteGroups.sglob.DzType.Rt.typeLab_of_local_m_alt}
\end{theorem}

\begin{proof}
\emph{Reduction to a local statement.} The norm data provide: the label truth of $D$; the representation parameters, i.e.\ the
local parameters (for the local Langlands correspondence of $h_L$) of the components of $\mathrm{rep}(\sigma)$ at the type places
are the local row parameters; the row depth, i.e.\ for all $\sigma,k,v$ a split place $s$ above $v$ and an irreducible
representation $\tau$ of $\GL_r(E_s)$ with the local row parameter of $\sigma$ and with non-zero vectors fixed by the principal
congruence subgroup $K_1$; and, at every $(k,v)$, an isomorphism $\kappa$ of the residue field at $v$ with $\overline{\F}_{g_k}$
coupling the tame parameter of the type at $(k,v)$ with the generator $\zeta_k$ of the group data, which equals the generator
$\mathrm{gen}_k$ of the label datum by the generator pin. Global Artin reciprocity under $h_L$ gives the large depth-zero inertia
condition (\lref{UpperTheory.Lit.Dz.global_artin_reciprocity_trivial}). The theorem therefore follows from the local
statement: under Green's and Zelevinsky's theorems and the large depth-zero inertia condition, the label truth, the
representation parameters, the row depth, the label tie and the tame coupling with the generators $\mathrm{gen}_k$ imply
$\mathrm{TL}(D,\mathrm{rep})$. We prove it.

\emph{The local statement.} Fix identifications of the completions at the type places
(\lref{GlobalTheory.T.nonempty_typePlaceCompletionIdent}). There is a \emph{class map} $\mathrm{rc}$, assigning to every
irreducible smooth representation at a type place a local parameter compatibly with the correspondence of $h_L$, such that an
irreducible representation with non-zero $K_1$-fixed vectors is determined up to isomorphism by its value under $\mathrm{rc}$
(\lref{UpperTheory.One.M.exists_isRc_isClassMapD_alt}). By the representation parameters and the row depth there are
irreducible components $\mathrm{cp}(\sigma,k,v)$ of $\mathrm{rep}(\sigma)$ at the type places with $\mathrm{rc}(\mathrm{cp}(\sigma,k,v))$
equal to the local row parameter (\lref{GlobalTheory.T.exists_comp_irrSm_locR_eq_rc}). Green's theorem, the compatibility of
tame generic supports with the class map (\lref{GlobalTheory.T.tame_gen_support_compat_of_is_rc}, from Zelevinsky's theorem),
the cuspidality of the depth-zero supercuspidal modules (\lref{GlobalTheory.T.depthZeroSupercuspidal_cuspidalModule}) and
the tame parameter bound obtained from the tame coupling and the large depth-zero inertia
(\lref{DiscreteGroups.sglob.DzType.Rt.tame_param_le_of_tame_coupled}) give the parameter bound for the labels
(\lref{DiscreteGroups.sglob.DzType.Rt.labParamLe_of}).

Next, every component $\mathrm{cp}(\sigma,k,v)$ has non-zero $K_1$-fixed vectors. Indeed, let $s,\tau$ be given by the row depth.
The pull-back $\tau'$ of $\tau$ to the rational local field has non-zero $K_1$-fixed vectors
(\lref{GlobalTheory.LabPin.principalCongruence_fixed_vector_rational_local_field}), and
$\mathrm{rc}(\tau')=\mathrm{rc}(\mathrm{cp}(\sigma,k,v))$, because both are the local row parameter of $\sigma$ and the map from
classes to parameters is injective. By the property of the class map, $\tau'\cong\mathrm{cp}(\sigma,k,v)$, and an isomorphism carries
non-zero $K_1$-fixed vectors to non-zero $K_1$-fixed vectors
(\lref{GlobalTheory.LabPin.principalCongruenceFixedVectors_ne_bot_of_iso}).

Finally, by the label equation of the row side, $\mathrm{lab}_k(\sigma,v)$ is the item attached to
$\mathrm{rc}(\mathrm{cp}(\sigma,k,v))$. Under Green's theorem, the label truth, the label tie and the parameter bound, a component with
non-zero $K_1$-fixed vectors has a label as item, and the maximal compact subgroup acts on its $K_1$-fixed vectors through the
series type of this label (\lref{DiscreteGroups.sglob.DzType.Rt.labLP_isLabel_and_actK_fix}); hence $\mathrm{rep}(\sigma)$ is fixed
by the series idempotent at every split place above $v$ (\lref{GlobalTheory.T.depthZeroIdentity_finiteType_of_typePrimeComponent}).
\end{proof}

For the label ties we use the data of a saturation $s$ of a point (\cref{pf:open0:lab}): a Hida setting $X$, a good operator $G$, a row
map $j$ from the rows to the square-integrable functions and the isometry $\iota$ of the spectral isomorphism; the real datum $R$
carries a theta map $\theta$, a multiplicity space $\mathrm{Mult}(\theta)$ for the level subgroup $J_R$ and the character $\chi_R$,
and an evaluation $e$ of the spectral vectors. For a row $\sigma$ put $e_\sigma=e(e_E(\mathrm{onb}_\sigma))$ (a multilinear function of
a vector $v$ and a group element $y$), and let $b_\sigma=\iota(s_{EE}(\mathrm{onb}_\sigma))$ be the base row of $\sigma$. The label
conditions of $B_{\mathrm{lab}}$ (\cref{pf:open0:def:twelve}) are: the labels of all finite types are labels, the label tuples are
good, and $b_\sigma$ lies in the type space $V_T(\mathrm{lab}\,\sigma)$ of its label tuple.

\begin{theorem}[Label ties]\label{pf:open0:irred}
Let $D$ be a datum with $F$ the maximal real subfield of $E$, $s$ a saturation over a real datum $R$ and a face $F$ with the standard
$\Psi$, and $\mathrm{rep}$ a type representation such that for every row $\sigma$ the vector $e_\sigma$ lies in $\mathrm{Mult}(\theta)$
and in its isotypic component of type $\mathrm{rep}(\mathrm{ft}\,\sigma)$. If $\mathrm{TL}(D,\mathrm{rep})$ holds and all finite types of $D$
are good, then the automorphic instance of $X$ matches $D$, the row dictionary tie holds, and the label conditions hold; that is,
$B_{\mathrm{lab}}$ holds for $(X,G,\iota)$.
\formallabel{GlobalTheory.T.constructedSpectralObjects_irreducibilityConditions}\ \formallabel{GlobalTheory.Sat.ties_at}\
\formallabel{GlobalTheory.T.realisationLabels_normalised_of_rowsTypes_levels}\ \formallabel{GlobalTheory.LabPinAt.labTie_at}\
\formallabel{GlobalTheory.T.labels_fixed_by_tupleOperator_of_typeSeries}\ \formallabel{GlobalTheory.LabPinTop.labTie_of_series}
\end{theorem}

\begin{proof}
The automorphic tie and the row dictionary tie are components of the spectral matching of $s$. The labels are labels by
$\mathrm{TL}(D,\mathrm{rep})$, and the label tuples are good since all finite types are good. It remains to show
$b_\sigma\in V_T(\mathrm{lab}\,\sigma)$ for a row $\sigma$.

\emph{Invariance of $e_\sigma$.} Vectors of the multiplicity space transform on the right under $J_R$ by $\chi_R$ and are left
invariant under the rational points (\lref{GlobalTheory.LabPinAt.mult_lev_left}). For $k$ in the level subgroup $K_g$ of $X$,
the level properties of $X$ give $\chi_2(k)=1$ (\lref{Formalism.L2.SG.Closed.level_subgroup_properties}), and the level match of $s$
gives $\kappa(k)\in J_R$ and $\chi_R(\kappa(k))=\overline{\chi_2(k)}$. Hence $e_\sigma(v,y\kappa(k))=e_\sigma(v,y)$
for $k\in K_g$, and $e_\sigma(v,\gamma y)=e_\sigma(v,y)$ for rational $\gamma$.

\emph{The series identity.} There are a reduction homomorphism $\mathrm{red}$ from the maximal compact subgroup to the finite group
$G_0$ of the level and a section $\mathrm{sec}$ of it, compatible with the standard frame
(\lref{GlobalTheory.T.exists_maximalCompact_reduction_standardFrame}). Let $(c_\gamma)_{\gamma\in G_0}$ be the coefficients of
the series idempotent of the label tuple of $\sigma$. For a function $\varphi$ in the isotypic component of type
$\mathrm{rep}(\mathrm{ft}\,\sigma)$ which is invariant under $K_g$,
$\sum_\gamma\overline{c_\gamma}\,\varphi(\cdot\,\kappa(\mathrm{sec}\,\gamma))=\varphi$
(\lref{GlobalTheory.T.sum_seriesIdempotentCoeff_rightTranslation} and \lref{GlobalTheory.T.fixedVectorIdentity_isotypic_conjugateType};
this uses $\mathrm{TL}(D,\mathrm{rep})$). Applied to $e_\sigma$ this gives
\begin{equation}\label{pf:open0:eq:series}
\sum_{\gamma\in G_0}\overline{c_\gamma}\,e_\sigma\bigl(v,y\,\kappa(\mathrm{sec}\,\gamma)\bigr)=e_\sigma(v,y)\qquad\text{for all }v,y.
\end{equation}

\emph{Transfer to the base row.} By the row pin, $b_\sigma$ is the complex conjugate of $j(\mathrm{onb}_\sigma)$, and $j(\mathrm{onb}_\sigma)$
is almost everywhere the function $x\mapsto a\,e_\sigma(v_0,\iota_X^{-1}(x)t)$ for some $v_0$, $a$ and a fixed element $t$. Thus
$b_\sigma(x)=\Phi(\iota_X^{-1}(x)t)$ almost everywhere, with $\Phi(y)=\overline{a\,e_\sigma(v_0,y)}$, which is left invariant under the
rational points. For such $\Phi$, right translation of $b_\sigma$ by $\mathrm{sec}\,\gamma$ is almost everywhere the function
$x\mapsto\Phi(\iota_X^{-1}(x)t\,\kappa(\mathrm{sec}\,\gamma))$ (\lref{GlobalTheory.LabPinTop.R_translate_ae}). Conjugating
\eqref{pf:open0:eq:series} with $v=v_0$ gives $\sum_\gamma c_\gamma\,R(\mathrm{sec}\,\gamma)b_\sigma=b_\sigma$. A vector fixed by the operator
$\sum_\gamma c_\gamma R(\mathrm{sec}\,\gamma)$ of a label tuple all of whose labels are labels lies in the type space of the tuple
(\lref{GlobalTheory.T.mem_rowsFixedByLabelTupleOp_of_sum_fixed}). Hence $b_\sigma\in V_T(\mathrm{lab}\,\sigma)$.
\end{proof}

\subsubsection{Isotypic components and the identification of \texorpdfstring{$\kappa$}{kappa}}\label{sec:pf:open0:kappa}

\begin{proposition}[Isotypic components]\label{pf:open0:isot}
Assume Harish-Chandra finiteness and Green's theorem. Let $t$ be row data of a twisted point (reading data with a good real datum,
pins and intersection data) for which the twisted multiplicities agree. Then for every row $\sigma$ and every $0\le i<m$ the
isotypic component $\mathrm{isot}_{i+1}(\mathrm{ft}\,\sigma)$ of the $(i+1)$-st row space of type $\mathrm{ft}\,\sigma$ is non-zero.
\formallabel{GlobalTheory.Fc.KapLinked.isot_ne_bot_of_rows}\ \formallabel{GlobalTheory.KapIsotRead.isot_ne_bot_read}
\end{proposition}

\begin{proof}
Green's theorem gives the idempotent tie for the types of the point
(\lref{GlobalTheory.DimPartsCompare.tameTie_of_typesTrue} and \lref{GlobalTheory.DimPartsCompare.TameTie.idemTie}); hence the coefficients of
the good idempotent are invariant under inversion, and the good operator is stable
(\lref{GlobalTheory.Mem.R3.GoodDual.goodStable_of_coeff}). The weak characterisation of the real datum gives the core
conditions on the level and the group (\lref{GlobalTheory.CoreOfWeak.core_of_weak}), from which the product of the character
idempotent and the level idempotent at the base is non-zero (\lref{GlobalTheory.GoodTypeGeneric.chiIdem_mul_lvlIdem_base_ne_zero}).
Under Harish-Chandra finiteness, the agreement of the multiplicities, the stability and this non-vanishing, the multilinear model of
the $i$-th row space contains a non-zero vector $w$ in its isotypic part of type $\mathrm{ft}\,\sigma$
(\lref{GlobalTheory.KapIsot.exists_mem_isotL}).

Let $\ell$ be the linear isomorphism from the multilinear model to the ambient space of the lift of the reading data. The vector
$\ell(w)$ lies in the good part, so $y=j_{M}^{-1}(\ell(w))\in M_{i+1}$ is defined, where $j_M$ is the isomorphism of the lift onto the good
part. By construction $j_M(y)$ lies in the part of type $\mathrm{ft}\,\sigma$, so by the lift tie
$y\in\mathrm{isot}_{i+1}(\mathrm{ft}\,\sigma)$. If this component were zero, then $y=0$, hence $\ell(w)=0$ and $w=0$, since $\ell$ is
injective; a contradiction.
\end{proof}

For a row side $S$ of $D$ on a finite-dimensional Hilbert space with operator $\iota\iota^*$, an orthonormal basis $(b_\sigma)_{\sigma\le n}$
with $n>0$ and types $\mathrm{ty}(\sigma)$, put $P(\sigma)=\prod_{j=1}^m\lambda'_{\mu_0,j}(\mathrm{ty}\,\sigma)$ and let $\kappa(S,\mathrm{ty})$ be
$\rho^{1/n}$ for $\rho=\Re\det(\iota\iota^*)/\prod_\sigma P(\sigma)$ if $\rho>0$. The \emph{sharp constant} $\kappa^\sharp(D)$ is
$\kappa(S^\sharp,\mathrm{ft})$ for the sharp row side $S^\sharp$ of $D$ with the basis $\mathrm{onb}$.

\begin{lemma}\label{pf:open0:kapofeigen}
If $P(\sigma)>0$ and $\iota\iota^*b_\sigma=\kappa\,P(\sigma)\,b_\sigma$ for all $\sigma$, with $\kappa>0$, then $\kappa(S,\mathrm{ty})=\kappa$.
\formallabel{GlobalTheory.KapRec.kapOf_of_eigen}
\end{lemma}

\begin{proof}
In the basis $(b_\sigma)$ the operator $\iota\iota^*$ is diagonal with entries $\kappa P(\sigma)$, so
$\det(\iota\iota^*)=\kappa^n\prod_\sigma P(\sigma)$ and $\rho=\kappa^n>0$. Hence $\kappa(S,\mathrm{ty})=(\kappa^n)^{1/n}=\kappa$.
\end{proof}

\begin{theorem}[The constant $\kappa$]\label{pf:open0:kappa}
Assume the discrete series cover, Harish-Chandra modules, Casselman--Wallach, Harish-Chandra finiteness and rapid decay. Let
$(R,F)$ be the real datum and face of reading data of a twisted point with pins and realisation conditions (which include
Theorem~B with a positive scalar $\kappa$), satisfying the $\lambda$-rows, and assume $\mathrm{isot}_{i+1}(\mathrm{ft}\,\sigma)\ne0$ for all
$\sigma$ and $i<m$. Let $X_d$ be a divisor datum with $\kappa_R(e)=\kappa^\sharp(D)$ for every complex embedding $e$ of $\mathit{LL}$.
Then $\kappa_R(e)$ equals the $\kappa$ of the row side of $(R,F)$ at the place of $e$, for every $e$.
\formallabel{GlobalTheory.T.divisorEmbedding_kappa_eq_rowsKappa_at_infinitePlace}\
\formallabel{GlobalTheory.KapRow.kapR_eq_rows_kappa}\
\formallabel{GlobalTheory.T.kappaConstant_sharpLevel_eq_dataKappaConstant}\ \formallabel{GlobalTheory.KapRec.kapSharp_eq_kap}
\end{theorem}

\begin{proof}
The $\kappa$ of the row side is the scalar $\kappa>0$ of Theorem~B at every place, and $\kappa_R(e)=\kappa^\sharp(D)$; so it suffices to
show $\kappa^\sharp(D)=\kappa$, by \cref{pf:open0:kapofeigen} for $S^\sharp$ and the basis $\mathrm{onb}$. The number $n$ of rows is positive
(\lref{GlobalTheory.Mem.R3.drows_pos}); every factor $\lambda'_{\mu_0,j}(\mathrm{ft}\,\sigma)$ is positive, using the non-vanishing of the
isotypic components (\lref{GlobalTheory.Mem.R3.LamPos.lambdaPrime_basis_row_type_pos_at_reference}); and $\iota\iota^*$ of the sharp
row side is the composite of the projection of the rows and its adjoint (\lref{GlobalTheory.KapRec.iiStar_apply_eq}), which has
$\mathrm{onb}_\sigma$ as eigenvector with eigenvalue $\kappa P(\sigma)$
(\lref{GlobalTheory.T.iotaProjectionAdjoint_basisRow_eigenvalue}: the basis rows lie in the type part of the definite family, on
which Theorem~B gives $\Theta\Theta^*=\kappa\prod_j\lambda'_j$, and the composite intertwiner is the theta lift composed with the good
projector).
\end{proof}

\subsubsection{Division compatibility}\label{sec:pf:open0:divc}

\begin{proposition}[Division compatibility]\label{pf:open0:divcompat}
Assume the discrete series cover, Harish-Chandra modules, Casselman--Wallach, Harish-Chandra finiteness and rapid decay. Let
$(R,F,Y,Z)$ be reading data of a twisted point with pins and realisation conditions, $C$ columns, $s_{EE}$ a conjugate-linear map and
$X_d$ a divisor datum, and assume the group tie (if the cardinalities and indices of the level groups agree, the group data agree).
\begin{enumerate}[label=(\alph*)]
\item The division definitions hold: ranks, group, tuples, the normalisation at $\mu_1$, the sharp and base definitions of the row
side, the metrics and lattices, the theta map and the columns.
\item If moreover the $M_j$-tie and the identification of the indices hold, the pairings satisfy the law \eqref{pf:open0:eq:law}, and the
base vanishing of \cref{pf:open0:btv} holds, then the division rows $\mathcal D_3$ hold.
\end{enumerate}
\formallabel{GlobalTheory.T.divisionCompatible_at_readIdentifications}\ \formallabel{GlobalTheory.divTied3_read}
\end{proposition}

\begin{proof}
(a) The pins give the cardinalities, indices and ranks of the level groups and the base values of the tuples
(\lref{GlobalTheory.Mem.R3.pins_card_index_ranks}); with the group tie they give the ranks and the group condition. The tuples, the
metrics and lattices and the columns are components of the reading data and of $C$; the sharp and base definitions of the row side hold
at both levels (\lref{GlobalTheory.T.rowSide_definition_both_levels}); and the theta map of the lifted cusp form space is the theta map
of the data (\lref{GlobalTheory.T.liftCuspFormSpace_thetaMap_eq_data}).

(b) The division rows consist of six conditions. The row lattices and metrics agree at the sharp and the base level
(\lref{GlobalTheory.T.rowLattices_metrics_agree_of_realGroupTheorems}); the zero subspace of the division rows is spanned by lifts of
the zero subspace of the data (\lref{GlobalTheory.T.zeroSubspace_spanned_by_dataZeroSubspace_lifts}); the theta tie is the second part
of the theta map identity of (a); the numbers of the rows are the concrete numbers at the base level, given the $M_j$-tie and the
identification of the indices (\lref{GlobalTheory.T.numbers_eq_concreteNumbers_at_baseLevel}); and the column pairings equal the
pairing $B_k$ of the divisor datum, given the law and the base vanishing (\lref{GlobalTheory.T.rationalPairing_eq_columnPairing}).
\end{proof}

\subsubsection{Inputs of the diagonal row theorem}\label{sec:pf:open0:diaginputs}

\paragraph{Identity cells.}
Let $X$ be a Hida setting for a datum $D$, $Y$ a full system and $u$ a word. In the archimedean group, let $N^U_\infty$ be the
unipotent subgroup and $H_\infty$ the reductive subgroup attached to $u$, with Haar measures $\mu_N$, $\mu_H$; let $\pi_u$ be the
representation of $u$ on the finite-dimensional space $V_\pi$, $\psi_\infty$ the character and $W_u$ the vector-valued function of
$u$, so that $W_u(sx)=\psi_\infty(s)\pi_u(s)W_u(x)$. For a term $t$ of the orbital expansion of the pair $(u,u)$ let $\gamma_t$ be the
archimedean component of its point, $R_\infty(t)$ its archimedean stabiliser and $\Omega_t$ its fundamental set; the \emph{orbital
operator} of $t$ is
\[
O(t)\xi=\int_{\Omega_t}\bigl\langle W_u(\gamma_t^{-1}x),\xi\bigr\rangle\,W_u(x)\,d\mu_\infty(x)\qquad(\xi\in V_\pi).
\]

\begin{lemma}\label{pf:open0:oot}
Let $t$ be a term with $\gamma_t^{-1}\in N^U_\infty$ such that $R_\infty(t)$ is a countable subgroup of $N^U_\infty$ with fundamental set
$\Omega_t$, let $\pi_u$ be irreducible, and let $\mu_H(H_\infty)$ and $\mu_N(\Omega_t)$ be finite. Then
\[
O(t)=\overline{\psi_\infty(\gamma_t^{-1})}\,\frac{\mu_N(\Omega_t)\,\mu_H(H_\infty)}{\dim V_\pi}\,\mathrm{id}_{V_\pi}.
\]
\formallabel{Formalism.L2.SG.TermOp.IdBlocks.oot_eq}
\end{lemma}

\begin{proof}
Since $\pi_u$ is trivial on $N^U_\infty$, $W_u(\gamma_t^{-1}x)=\psi_\infty(\gamma_t^{-1})W_u(x)$; as the inner product is
conjugate-linear in its first argument, $O(t)=\overline{\psi_\infty(\gamma_t^{-1})}\int_{\Omega_t}\langle W_u(x),\cdot\rangle W_u(x)\,d\mu_\infty$.
The remaining integral is $\mu_N(\Omega_t)\mu_H(H_\infty)/\dim V_\pi$ times the identity by the orthogonality relation of the development
(\lref{Formalism.L2.SG.TermOp.IdBlocks.oo_eq}: under the disintegration of $\mu_\infty$ over the product of $N^U_\infty$ and $H_\infty$, the
equivariance of $W_u$, the normalisation of $W_u$, the unitarity and irreducibility of $\pi_u$ on $H_\infty$ and the invariance of $\mu_N$
under the countable group $R_\infty(t)$, Schur orthogonality evaluates the integral over a fundamental set).
\end{proof}

The column space $\mathrm{CM}(u)$ of the column system of $Y$ at $u$ is a finite orthogonal sum of copies of $V_\pi$ indexed by
classes $\rho$, and the operator $a(t)$ of a term $t$ on $\mathrm{CM}(u)$ has blocks $a(t)_{\rho\rho'}\in\End(V_\pi)$. A block vanishes
unless $t$ \emph{meets} $(\rho,\rho')$, in which case
$a(t)_{\rho\rho'}=(c_0^2)^{-1/2}\,\vol(K_2)\,\chi_2(k_t)\,O(t)$ (\lref{Formalism.L2.SG.TermOp.IdBlocks.blk_of}), with $k_t$ a level element
attached to $t$ and the normalising constant $c_0=\vol(K_2)\,\mu_N(F_N)\,\mu_H(H_\infty)/\dim V_\pi$ of $u$, where $F_N$ is the fundamental
set of $N^U_\infty$ used for the expansion.

\begin{lemma}[Identity cells]\label{pf:open0:idcells}
Let $D$ be a datum, $X$ a Hida setting with normalised level subgroups, $Y$ a full system and $u$ a word for which fundamental sets
exist and the orbital integrals of the pair $(u,u)$ are integrable. Then the operators $a(t)$ of the identity terms $t$ of the orbital
expansion of $(u,u)$ have unconditional sum $\mathrm{id}_{\mathrm{CM}(u)}$.
\formallabel{Formalism.L2.SG.Closed.IdCellDiagC_holds}\ \formallabel{Formalism.L2.SG.Pf.idCellDiagC_of_steps3}\
\formallabel{Formalism.L2.SG.Pf.idMeetC_of_steps3}\ \formallabel{Formalism.L2.SG.TermOp.IdBlocks.idCellDiagC_of_steps}\
\formallabel{Formalism.L2.SG.TermOp.KExpCut.idCellDiagC_of}\ \formallabel{Formalism.L2.SG.TermOp.IdBlocks.idDiagBlocksC}\
\formallabel{Formalism.L2.SG.TermOp.IdBlocks.idDiagBlocks_of}\ \formallabel{Formalism.L2.SG.TermOp.KExpCut.idCellDiag_of}
\end{lemma}

\begin{proof}
Since $\mathrm{CM}(u)$ is the finite orthogonal sum of the blocks and an operator is determined continuously by its blocks, it suffices
to show that for all classes $\rho,\rho'$ the family $(a(t)_{\rho\rho'})_t$ over the identity terms has sum $\mathrm{id}_{V_\pi}$ if
$\rho=\rho'$ and $0$ otherwise. We use the
following properties of identity terms, valid under the hypotheses since the level subgroup at $u$ is neat
(\lref{Formalism.L2.SG.Closed.NeatC_holds}):
\begin{enumerate}[label=(\roman*)]
\item an identity term meeting $(\rho,\rho')$ has $\rho=\rho'$, and for every $\rho$ exactly one identity term meets $(\rho,\rho)$
(\lref{Formalism.L2.SG.Pf.idMeetOneC_holds});
\item if an identity term $t$ meets $(\rho,\rho)$, then $\gamma_t^{-1}\in N^U_\infty$ and $\chi_2(k_t)\,\overline{\psi_\infty(\gamma_t^{-1})}=1$
(\lref{Formalism.L2.SG.Pf.idMeetPhaseC_holds});
\item for such $t$, $R_\infty(t)$ is a countable subgroup of $N^U_\infty$ (\lref{Formalism.L2.SG.Pf.idMeetRO_of}, from the inclusion of
the archimedean components of rational links into $N^U_\infty$, \lref{Formalism.L2.SG.Pf.nLamNinf_holds});
\item for such $t$, $\mu_N(\Omega_t)=\mu_N(F_N)$ (\lref{Formalism.L2.SG.Pf.idMeetVolC_holds}).
\end{enumerate}
Moreover $\pi_u$ is irreducible (\lref{Formalism.L2.SG.Pf.piIrrC_holds}), $\vol(K_2)>0$ by the normalisation of the level subgroups,
and $\mu_H(H_\infty)$ and $\mu_N(F_N)$ are positive and finite (\lref{Formalism.L2.SG.Closed.VolPosC_holds}).

If $V_\pi=0$ all operators vanish and there is nothing to prove. Otherwise, for $\rho\ne\rho'$ no identity term meets $(\rho,\rho')$ by
(i), so all blocks vanish. For $\rho=\rho'$ let $t_0$ be the identity term meeting $(\rho,\rho)$; by (i) the blocks of all other
identity terms vanish, so the sum is $a(t_0)_{\rho\rho}$. By \cref{pf:open0:oot}, applicable by (ii), (iii) and the finiteness of the
measures, and by (iv),
\[
a(t_0)_{\rho\rho}=(c_0^2)^{-1/2}\vol(K_2)\,\chi_2(k_{t_0})\overline{\psi_\infty(\gamma_{t_0}^{-1})}\,
\frac{\mu_N(F_N)\,\mu_H(H_\infty)}{\dim V_\pi}\,\mathrm{id}=c_0^{-1}c_0\,\mathrm{id}=\mathrm{id},
\]
using (ii) and $(c_0^2)^{-1/2}=c_0^{-1}$, as $c_0>0$.
\end{proof}

\paragraph{The pair amplitude bound.}
For concretely tied $(X,Y)$ and a word $u$, the \emph{type condition} $\mathrm{Ty}(C_H,C_s,C,C_h,C_p)$ at $u$ is the type bound of the
development cited in \cref{pf:open0:deg}.

\begin{lemma}\label{pf:open0:pg2}
Under the $K_2$-hypotheses, $\openzeroPg(Y)\le2$.
\formallabel{GlobalTheory.Core.Deg.pgC_le_two_at_record}\ \formallabel{Formalism.L2.SG.KappaOffDiagonal.PgLe.one_factor}
\end{lemma}

\begin{proof}
The typed reduction of $Y$ consists of the places above $g_0$ and $g_1$, so $\openzeroPg(Y)=(1+3/g_0)^f(1+3/g_1)^f$
(\lref{Formalism.L2.SG.Pf.SatGram.tt_eq'}). The $K_2$-hypotheses give $(kb)^8\le g_i$, $f\le k$, $k\ge1$ and $b\ge2$. Hence
$g_i\ge256k^8$ and $(1+3/g_i)^f\le e^{3f/g_i}\le e^{3k/(256k^8)}\le e^{1/4}$, so that $\openzeroPg(Y)\le e^{1/2}\le2$.
\end{proof}

\begin{proposition}[Pair amplitude bound]\label{pf:open0:pairsamp}
For all $C_H,C_h,C_p\ge0$ and $C\ge2(C_H+1)$ there is $C_{l,0}$ such that for all constants with $C_l\ge C_{l,0}$, every datum $D$ with
size data $L_2$, every concretely tied pair $(X,Y)$ and every word $u$ satisfying the $K_2$-hypotheses and $\mathrm{Ty}(C_H,C_p+4,C,C_h,C_p)$
at $u$, the family $(\|\mathrm{vec}\,t\|\,\|\mathrm{vec}\,t'\|)_{(t,t')\in P(\openzeroDeg)}$ of $\mathcal T(u,u)$ is summable and
\[
\openzeroPg(Y)\,\Gamma_1\sum_{(t,t')\in P(\openzeroDeg)}\|\mathrm{vec}\,t\|\,\|\mathrm{vec}\,t'\|\le d_u\,e^{-\openzerokfk}.
\]
\formallabel{GlobalTheory.T.pairsAmplitudeBound_of_typeBound}\ \formallabel{GlobalTheory.Core.Deg.pairsDegOfType_holds}
\end{proposition}

\begin{proof}
Put $C_c=21+4|C_H|$, let $C_f\ge1$ be the explicit constant formed from $(C,C_h,C_c)$ in the amplitude estimates of the development, and
let $C_{l,0}=\max(1,K(C_f))$, where $K(C_f)$ is the constant of the final estimate below. Fix data as in the statement.

\emph{Inputs.} (1) Every relevant term $t$ has a positive invariant $\Delta(t)$, bounded by an explicit factor depending on $b,C,m,f$ times the square of an explicit size of $t$ (\lref{Formalism.L2.SG.KappaOffDiagonal.delta_pos_le_height_sq}). (2) The label sums of the record are
bounded (\lref{GlobalTheory.Core.Deg.sum_labels_one_record}). (3) $\openzeroPg(Y)\le2$ (\cref{pf:open0:pg2}). (4) By the dimension
formula, $\Gamma_1\ge0$ and $\Gamma_1\,m(u)\,d_\pi\le d_u$, where $m(u)\ge0$ is the mass of the rows of $u$ and $d_\pi=\dim V_\pi$
(\lref{kappa_off_diagonal_dim_form_C}). (5) The positivity of the $J$-factors holds (\lref{Formalism.L2.SG.KappaOffDiagonal.JPosC_holds}).
(6) Fundamental sets exist and the orbital integrals are integrable for concretely tied pairs
(\lref{Formalism.L2.SG.Pf.fundExistsC_holds} and \lref{Formalism.L2.SG.Closed.OOIntAllC_outright}), so degenerate terms which are not
relevant have vector $0$ (\lref{Formalism.L2.SG.KappaOffDiagonal.IrrVecC_holds}).

\emph{Partial sums.} Let $A_1(j)\ge0$ ($1\le j\le r_S$) be the explicit rank contributions of the development and
$B=d_\pi\,m(u)\bigl(\sum_{j=1}^{r_S}A_1(j)\bigr)^2$. From (1), (2), (5), (6) and the type condition, every finite partial sum of the family
over $P(\openzeroDeg)$ is at most $B$ (\lref{Formalism.L2.SG.KappaOffDiagonal.fin_bound}). As the family is non-negative, it is summable
with sum at most $B$.

\emph{The final estimate.} Each $A_1(j)$ is at most the explicit term $A(j)$ of the development
(\lref{Formalism.MainNotions.KappaEndpoints.rank_contribution_le}), so $\bigl(\sum_jA_1(j)\bigr)^2\le B_2$ for the explicit
$B_2=\bigl(\sum_jA(j)\bigr)^2$. Since $C_l\ge K(C_f)$ and the $K_2$-hypotheses give $C_lm^6L\le k\log k$ with $L\ge0$, we have
$K(C_f)m^6L\le k\log k$, and then $2d_uB_2\le d_ue^{-\openzerokfk}$ (\lref{Formalism.MainNotions.KappaEndpoints.final_bound_closes}).
Hence, by (3) and (4),
\[
\openzeroPg(Y)\,\Gamma_1\sum_{P(\openzeroDeg)}\|\mathrm{vec}\,t\|\,\|\mathrm{vec}\,t'\|\le\openzeroPg(Y)\,(\Gamma_1m(u)d_\pi)\Bigl(\sum_jA_1(j)\Bigr)^2
\le2\,d_u\,B_2\le d_u\,e^{-\openzerokfk}.
\qedhere
\]
\end{proof}

\paragraph{Inputs of the collision bound.}
For a datum $D$ with size data $L_2$ the \emph{amplitude bounds} $\mathrm{Amp}(C',C_l,C_0;C_{\mathrm{pr}},C,C_2,C_h,\delta)$ are the three
estimates for the amplitudes at the size data used in \cref{pf:open0:coll}: an $L^2$-bound and an absolute bound for products of
amplitudes, and a decay bound with exponent $\delta$.

\begin{proposition}[Amplitude bounds]\label{pf:open0:ampon}
Let $c\in\R$. For all $C_g,c_b,C_u,c_T,C_T$ with $c_b,c_T>0$ there are $C,C_h\ge0$ and $C_2$ such that for all $C_{\mathrm{pr}}$ and
$\delta>0$ there is a threshold $t$ such that $\mathrm{Amp}(C',C_l,c;C_{\mathrm{pr}},C,C_2,C_h,\delta)$ holds for all $C',C_l\ge t$ and every
datum $D$ with size data $L_2$ such that $C_\Gamma(L_2)\ge1$, the ties of $D$, the mid-volume formula and the trace frame hold.
\formallabel{GlobalTheory.T.rowAmplitudeBoundsUniform_singleton_width}\ \formallabel{GlobalTheory.AmpRecordGap.ampOn_singleton}\
\formal{GlobalTheory.Core.Deg.AmpAt}
\end{proposition}

\begin{proof}
The amplitude theorem of the development (\lref{GlobalTheory.T.ampAt_of_admitted_bounds}) states: for constants
$C_g,C_{PT},C_{Dl}$, a Hilbert-90 constant $h_{90}\ge0$, $C_a\ge0$ and $C_S$ there are $C,C_h\ge0$ and $C_2$ such that for all $C_{\mathrm{pr}}$
and $\delta>0$ there is a threshold $t_0$ beyond which the amplitude bounds hold at every datum with $C_\Gamma\ge1$ satisfying eight
admissibility conditions (A1)--(A8). We take $h_{90}=1$ and obtain the constants and the conditions from the following inputs, each of
which provides, for given $C_{\mathrm{pr}}$, a threshold beyond which it holds at the data of the statement:
(A1) the wide-range bound for ample theta series, with a constant $C_{PT}$, which uses the trace frame
(\lref{GlobalTheory.T.ampleTheta_wideRange_large_constants});
(A2) the point bound for $D_l$, with a constant $C_{Dl}$ (\lref{GlobalTheory.CoreDlGap.ptDl_at_record});
(A3) the low-weight condition: under the standing hypotheses, the critical level is bounded below
(\lref{GlobalTheory.T.exists_threshold_criticalLevel_lower_bound}), using that the normalised formal degrees are Vandermonde
expressions (\lref{GlobalTheory.T.normalizedFormalDegree_eq_vandermonde_of_sizes}) and that the standing numerical data contain the
conditions $\mathrm{NW}$ (\lref{GlobalTheory.T.dataNW_of_standingNumericalData});
(A4) the bound for the real-place amplitude by a constant divided by the formal degree
(\lref{GlobalTheory.T.realPlaceAmplitude_norm_le_const_div_formalDegree});
(A5) the unit-shift bound for degree ratios with $h_{90}=1$: along the line of the word the ratio of formal degrees is a ratio of
Vandermonde expressions (\lref{GlobalTheory.T.formalDegree_ratio_eq_vandermonde_ratio_along_line}), which gives the bound
(\lref{Formalism.MainNotions.BridgeLemma₂.degree_ratio_one_of_vandermonde});
(A6) the absolute bounds at all parameters, with a constant $C_a\ge0$ (\lref{GlobalTheory.AbsRecGap.abs_all_at_record_gap});
(A7) the Schur bound, with a constant $C_S$ (\lref{GlobalTheory.Amp.schur_at_record_gap});
(A8) the Weyl condition, which holds for every datum (\lref{GlobalTheory.AmpRecordGap.weyl_at}).
Let $C,C_2,C_h$ be the constants of the amplitude theorem for $(C_g,C_{PT},C_{Dl},1,C_a,C_S)$. Given $C_{\mathrm{pr}}$ and $\delta>0$, let $t$ be
the maximum of $t_0$ and of the thresholds of (A1)--(A7) for $\max(C_{\mathrm{pr}},0)$. For $C',C_l\ge t$ and $C_0=c$ all eight conditions hold,
and the amplitude theorem gives the amplitude bounds.
\end{proof}

\begin{proposition}[Term and loop ties]\label{pf:open0:tieon}
Let $c\in\R$. For all $C_g,c_b,C_u,c_T,C_T$ there are $c_V$ and a threshold $t$ such that for all $C',C_l\ge t$, every datum $D$ with size data
$L_2$, every Hida setting $X$, every full system $Y$ with finite spaces and isometric identification of the spectral data, and every
word $u$, if the ties of $D$, $\mathrm{H}_0(c)$ and the concrete tie of $(X,Y)$ hold, then there are a term tie $T$ and a loop tie $Z$ at
$(u,u)$ satisfying the tie conditions and the volume bound with constant $c_V$.
\formallabel{GlobalTheory.Core.Deg.tieOn_singleton_uncond}
\end{proposition}

\begin{proof}
Let $C$ be the constant of the covolume bound for stabilisers of relevant terms
(\lref{GlobalTheory.T.stabiliserCovolume_relevantTerm_le_heightPow}: under the standing hypotheses and Wu data, the covolume of the
stabiliser of every relevant term of non-zero weight is at most an explicit power of its height, with constant $C$), and put $c_V=C$ and
$t=72$. For $C',C_l\ge72$ and data as in the statement, the term-loop theorem
(\lref{GlobalTheory.Core.Deg.exists_termLoop_of_standingHypotheses}, valid under the ties, the standing hypotheses and the concrete tie
for $C_l\ge72$; its proof combines the gap hypothesis, Wu decay, the uniqueness of Bessel models with large gap, the covolume formula for
unipotent integrals and the inequality $e_N\log(2v)\le m^2\log b$, which follows from $\log\mathrm{rd}_E\le L/C_l$, $C_l\ge72$ and
$e_N\le m^2/2$) provides Wu data, a term tie $T$ and a loop tie $Z$ satisfying the tie conditions, and the
identity of the volume $\vol_Z(t)$ of every term $t$ with the covolume of its stabiliser. The volume bound therefore follows from the
covolume bound.
\end{proof}

\begin{proposition}[Explicit size hypothesis]\label{pf:open0:xisizes}
Assume the Siegel mass ratio. For every $C\ge408\,e\pi$, every datum $D$ with size data $L_2$, all words $u,u'$ and all constants with
$C_l\ge1$ such that the count hypotheses for $(u,u')$ hold, the size hypothesis $\mathrm{XH}(C)$ of \cref{pf:open0:corank} holds for the size
function $\Xi$ of $(u,u')$, with $m=2r_S+2$ and $N_T=T_u^f$.
\formallabel{GlobalTheory.T.corankSizeHypothesis_explicitConstant_siegelMass}\ \formallabel{GlobalTheory.Core.Deg.xiSizes_explicit}
\end{proposition}

\begin{proof}
Since $C_l\ge1$, $8\max(1,1/C_l)\le8$, so the plane mass bound with parameter $8$ holds
(\lref{Formalism.MainNotions.BridgeLemma₂.plane_mass_bound_parameters}: from short vectors at every place, the count of plane classes
and the existence of a second orthogonal vector). The Siegel mass ratio gives the first lattice count
(\lref{Formalism.MainNotions.BridgeLemma.lattice_count_first_of_siegel}), and with the plane mass bound the first size bound with
constant $408$ (\lref{Formalism.MainNotions.BridgeLemma.first_bound_of_mass_ratio}). Since $e\ge1$ and $\pi\ge3$, $e\pi\ge1$ and
$408\max(1,e\pi)=408\,e\pi\le C$. Under the count hypotheses, the first size bound and this inequality give $\mathrm{XH}(C)$
(\lref{Formalism.MainNotions.BridgeLemma.xiHyp_of_stX}); finally $m=2r_S+2$ (\lref{Formalism.MainNotions.BridgeLemma.d_m_eq}).
\end{proof}

\paragraph{The generic corank-one row.}
The generic row $\mathrm{R}^{\mathrm{gen}}$ was stated before \cref{pf:open0:onerows}. The hypotheses of the following theorem are the
volume identities $\sfname{SWXS}$ of the Siegel--Weil formula and the step statements listed in the proof of
\cref{pf:open0:fourpoints}(b).

\begin{theorem}[Generic corank-one row]\label{pf:open0:gen1}
Assume the volume identities $\sfname{SWXS}$, the normal line, integrality and separation statements, the point and norm-one steps, the
tie of contexts, the step choice constants, the existence of step pins and the step Theorem~B. Then $\mathrm{R}^{\mathrm{gen}}$ holds.
\formallabel{GlobalTheory.Core.One.rowGen1P_of_siegelWeil_of_points}
\end{theorem}

\begin{proof}
\emph{Classes.} Within a row, corank one is an equivalence relation on the terms of non-zero weight
(\lref{GlobalTheory.Core.One.corankOne_isEquivalence_nonzeroGram_terms}); this gives the class step
(\lref{GlobalTheory.Core.One.stepClass1_of}) and the family data of the generic classes (\lref{GlobalTheory.Core.One.famGen1}).

\emph{The family step.} The local mass is at most nine times the generic mass (\lref{GlobalTheory.localMass_le_genericMass_nine});
with the leaf tie (\lref{GlobalTheory.HTieLeaf.main}), the unfolding identity for family row-weight sums
(\lref{GlobalTheory.SnFam.familyRowWeightSumUnfoldingIdentity_proof}), the family weights obtained from the volume identities
(\lref{GlobalTheory.WeightsFamPf.weightsFam_of}), the factorisation of family group sums into row-weight amplitudes
(\lref{GlobalTheory.FamilyNormRational.familyGroupSumFactorsRowWeightAmplitude}) and the existence of families
(\lref{GlobalTheory.SnFam.existFam}), the family data give the family step (\lref{GlobalTheory.Core.One.stepFam1_of}). The local
non-generic bound for families with constant $39$ (\lref{GlobalTheory.local_nongeneric_bound_record_family}) and the family step give
the local step (\lref{GlobalTheory.Core.One.stepLoc1_of_fam}).

\emph{Normals and vanishing.} The equivalence of corank one, the normal line, integrality and separation, the normal count
(\lref{Formalism.L2.SG.NGCut.NormalCount.NormCount_holds}) and the norm-one step give the normal step
(\lref{GlobalTheory.Core.One.stepNormals1_of}). The form relation of the development gives the form condition
(\lref{Formalism.L2.SG.Pf.formRelC_holds} and \lref{Formalism.L2.SG.Pf.formGC_of_rel}), which gives the term vanishing step
(\lref{Formalism.L2.SG.Pf.MainEqVanish.stepTermVanishG_of_form}).

\emph{Assembly.} The generic pairs of corank one are grouped by the rational points of the auxiliary space; for each point the group sum
is expressed through mass data whose family weights come from the volume identities, and a joint row bound over the points gives
(G1)$_{\gamma_L(C_B)}$. Formally, the norm chain (\lref{Formalism.FourierDecayAlt.P.norm_chain}), the class, normal, point, norm-one,
tie, choice, pin, vanishing and local steps and the step Theorem~B imply $\mathrm{R}^{\mathrm{gen}}$
(\lref{GlobalTheory.Core.One.rowGen1P_of_steps}).
\end{proof}

%==============================================================================
\subsection{Remarks on the formal statements}\label{sec:pf:open0:remarks}

\begin{remark}[Unused hypotheses]\label{pf:open0:rem:unused}
(a) Four of the twelve statements hold for every parameter and use no classical input: primitivity, decomposition
points, comparisons and the definite-side row (\cref{pf:open0:hpa,pf:open0:dcpt,pf:open0:tie,pf:open0:df}). The formal statements
of the comparison row and of the definite-side row nevertheless carry seven, resp.\ four, classical results as hypotheses,
and the former also a hypothesis on the extension; all are discarded.
(b) The formal statement of the join row also assumes the primitivity and decomposition-point rows, the existence of
Weil representations and of discrete series; these are not used, only the comparison row and five classical inputs. The
division row does not use the existence of Weil representations and of discrete series either.
(c) Of the fields of $P$ listed after \cref{pf:open0:main}, each is passed to at least one of the statements; whether each is
used inside the cited ingredients has not been checked here. The remaining fields of $P$ (Galois-theoretic, $p$-adic Hodge
theoretic and Bernstein--Zelevinsky inputs, among others) are not used by the theorem.
\end{remark}

\begin{remark}[Where the content lies]\label{pf:open0:rem:content}
A bundle is linked to a point only through a bundle witness, which identifies its real package with the real package of a
twisted input built at the point; the primitivity, decomposition-point and comparison statements are read off this
identification. The auxiliary slot of $\openzeroLinked$ is the trivial slot for all statements except goodness and labels,
which use the slot of points with good finite types, and the geometric row, which inherits the slot of the theorem on
class-sum divisors.
None of the twelve statements asserts that linked bundles exist; this is the subject of \cref{thm:member-bundles}. The analytic
content specific to this section is the diagonal row theorem (\cref{pf:open0:diagthm}), the input of the non-vanishing of periods
in \cref{sec:pf:open0a}; it is a diagonal companion of Theorem~S (\cref{thm:theorem-S}), proved with the same orbital expansion,
but for equal words, where the identity terms contribute the identity operator and only corank-one pairs, collisions and
degenerate terms remain.
\end{remark}

\begin{remark}[The parameter $c$]\label{pf:open0:rem:c}
The real number $c$ enters $\openzeropdep(h_L,c)$ only through the list pin of the extension $N_{\mathrm{dep}}$. No step of the proof
uses the list pin: every proposition of this section holds for all parameters $p(N,S,\Gamma)$ whose extension yields the law
extension (or for all parameters), so \cref{pf:open0:main} holds uniformly in $c$ for the trivial reason that the list pin is never
used. The constant $c$ in the definite-side row (the third argument of the bundles) is unrelated: there it is the size constant
$C_{\mathrm{sz}}$ of the family.
\end{remark}


